\documentclass[
  doc,
  floatsintext,
  longtable,
  a4paper,
  nolmodern,
  notxfonts,
  notimes,
  12pt,
  colorlinks=true,linkcolor=blue,citecolor=blue,urlcolor=blue]{apa7}

\usepackage{amsmath}
\usepackage{amssymb}

\geometry{inner=1in, outer=1in}
\fancyhfoffset[LE,RO]{0cm}


\usepackage[bidi=default]{babel}
\babelprovide[main,import]{spanish}


% get rid of language-specific shorthands (see #6817):
\let\LanguageShortHands\languageshorthands
\def\languageshorthands#1{}

\RequirePackage{longtable}
\RequirePackage{threeparttablex}

\makeatletter
\renewcommand{\paragraph}{\@startsection{paragraph}{4}{\parindent}%
	{0\baselineskip \@plus 0.2ex \@minus 0.2ex}%
	{-.5em}%
	{\normalfont\normalsize\bfseries\typesectitle}}

\renewcommand{\subparagraph}[1]{\@startsection{subparagraph}{5}{0.5em}%
	{0\baselineskip \@plus 0.2ex \@minus 0.2ex}%
	{-\z@\relax}%
	{\normalfont\normalsize\bfseries\itshape\hspace{\parindent}{#1}\textit{\addperi}}{\relax}}
\makeatother




\usepackage{longtable, booktabs, multirow, multicol, colortbl, hhline, caption, array, float, xpatch}
\usepackage{subcaption}


\renewcommand\thesubfigure{\Alph{subfigure}}
\setcounter{topnumber}{2}
\setcounter{bottomnumber}{2}
\setcounter{totalnumber}{4}
\renewcommand{\topfraction}{0.85}
\renewcommand{\bottomfraction}{0.85}
\renewcommand{\textfraction}{0.15}
\renewcommand{\floatpagefraction}{0.7}

\usepackage{tcolorbox}
\tcbuselibrary{listings,theorems, breakable, skins}
\usepackage{fontawesome5}

\definecolor{quarto-callout-color}{HTML}{909090}
\definecolor{quarto-callout-note-color}{HTML}{0758E5}
\definecolor{quarto-callout-important-color}{HTML}{CC1914}
\definecolor{quarto-callout-warning-color}{HTML}{EB9113}
\definecolor{quarto-callout-tip-color}{HTML}{00A047}
\definecolor{quarto-callout-caution-color}{HTML}{FC5300}
\definecolor{quarto-callout-color-frame}{HTML}{ACACAC}
\definecolor{quarto-callout-note-color-frame}{HTML}{4582EC}
\definecolor{quarto-callout-important-color-frame}{HTML}{D9534F}
\definecolor{quarto-callout-warning-color-frame}{HTML}{F0AD4E}
\definecolor{quarto-callout-tip-color-frame}{HTML}{02B875}
\definecolor{quarto-callout-caution-color-frame}{HTML}{FD7E14}

%\newlength\Oldarrayrulewidth
%\newlength\Oldtabcolsep


\usepackage{hyperref}



\usepackage{color}
\usepackage{fancyvrb}
\newcommand{\VerbBar}{|}
\newcommand{\VERB}{\Verb[commandchars=\\\{\}]}
\DefineVerbatimEnvironment{Highlighting}{Verbatim}{commandchars=\\\{\}}
% Add ',fontsize=\small' for more characters per line
\usepackage{framed}
\definecolor{shadecolor}{RGB}{241,243,245}
\newenvironment{Shaded}{\begin{snugshade}}{\end{snugshade}}
\newcommand{\AlertTok}[1]{\textcolor[rgb]{0.68,0.00,0.00}{#1}}
\newcommand{\AnnotationTok}[1]{\textcolor[rgb]{0.37,0.37,0.37}{#1}}
\newcommand{\AttributeTok}[1]{\textcolor[rgb]{0.40,0.45,0.13}{#1}}
\newcommand{\BaseNTok}[1]{\textcolor[rgb]{0.68,0.00,0.00}{#1}}
\newcommand{\BuiltInTok}[1]{\textcolor[rgb]{0.00,0.23,0.31}{#1}}
\newcommand{\CharTok}[1]{\textcolor[rgb]{0.13,0.47,0.30}{#1}}
\newcommand{\CommentTok}[1]{\textcolor[rgb]{0.37,0.37,0.37}{#1}}
\newcommand{\CommentVarTok}[1]{\textcolor[rgb]{0.37,0.37,0.37}{\textit{#1}}}
\newcommand{\ConstantTok}[1]{\textcolor[rgb]{0.56,0.35,0.01}{#1}}
\newcommand{\ControlFlowTok}[1]{\textcolor[rgb]{0.00,0.23,0.31}{\textbf{#1}}}
\newcommand{\DataTypeTok}[1]{\textcolor[rgb]{0.68,0.00,0.00}{#1}}
\newcommand{\DecValTok}[1]{\textcolor[rgb]{0.68,0.00,0.00}{#1}}
\newcommand{\DocumentationTok}[1]{\textcolor[rgb]{0.37,0.37,0.37}{\textit{#1}}}
\newcommand{\ErrorTok}[1]{\textcolor[rgb]{0.68,0.00,0.00}{#1}}
\newcommand{\ExtensionTok}[1]{\textcolor[rgb]{0.00,0.23,0.31}{#1}}
\newcommand{\FloatTok}[1]{\textcolor[rgb]{0.68,0.00,0.00}{#1}}
\newcommand{\FunctionTok}[1]{\textcolor[rgb]{0.28,0.35,0.67}{#1}}
\newcommand{\ImportTok}[1]{\textcolor[rgb]{0.00,0.46,0.62}{#1}}
\newcommand{\InformationTok}[1]{\textcolor[rgb]{0.37,0.37,0.37}{#1}}
\newcommand{\KeywordTok}[1]{\textcolor[rgb]{0.00,0.23,0.31}{\textbf{#1}}}
\newcommand{\NormalTok}[1]{\textcolor[rgb]{0.00,0.23,0.31}{#1}}
\newcommand{\OperatorTok}[1]{\textcolor[rgb]{0.37,0.37,0.37}{#1}}
\newcommand{\OtherTok}[1]{\textcolor[rgb]{0.00,0.23,0.31}{#1}}
\newcommand{\PreprocessorTok}[1]{\textcolor[rgb]{0.68,0.00,0.00}{#1}}
\newcommand{\RegionMarkerTok}[1]{\textcolor[rgb]{0.00,0.23,0.31}{#1}}
\newcommand{\SpecialCharTok}[1]{\textcolor[rgb]{0.37,0.37,0.37}{#1}}
\newcommand{\SpecialStringTok}[1]{\textcolor[rgb]{0.13,0.47,0.30}{#1}}
\newcommand{\StringTok}[1]{\textcolor[rgb]{0.13,0.47,0.30}{#1}}
\newcommand{\VariableTok}[1]{\textcolor[rgb]{0.07,0.07,0.07}{#1}}
\newcommand{\VerbatimStringTok}[1]{\textcolor[rgb]{0.13,0.47,0.30}{#1}}
\newcommand{\WarningTok}[1]{\textcolor[rgb]{0.37,0.37,0.37}{\textit{#1}}}

\providecommand{\tightlist}{%
  \setlength{\itemsep}{0pt}\setlength{\parskip}{0pt}}
\usepackage{longtable,booktabs,array}
\usepackage{calc} % for calculating minipage widths
% Correct order of tables after \paragraph or \subparagraph
\usepackage{etoolbox}
\makeatletter
\patchcmd\longtable{\par}{\if@noskipsec\mbox{}\fi\par}{}{}
\makeatother
% Allow footnotes in longtable head/foot
\IfFileExists{footnotehyper.sty}{\usepackage{footnotehyper}}{\usepackage{footnote}}
\makesavenoteenv{longtable}

\usepackage{graphicx}
\makeatletter
\newsavebox\pandoc@box
\newcommand*\pandocbounded[1]{% scales image to fit in text height/width
  \sbox\pandoc@box{#1}%
  \Gscale@div\@tempa{\textheight}{\dimexpr\ht\pandoc@box+\dp\pandoc@box\relax}%
  \Gscale@div\@tempb{\linewidth}{\wd\pandoc@box}%
  \ifdim\@tempb\p@<\@tempa\p@\let\@tempa\@tempb\fi% select the smaller of both
  \ifdim\@tempa\p@<\p@\scalebox{\@tempa}{\usebox\pandoc@box}%
  \else\usebox{\pandoc@box}%
  \fi%
}
% Set default figure placement to htbp
\def\fps@figure{htbp}
\makeatother







\usepackage{newtx}

\defaultfontfeatures{Scale=MatchLowercase}
\defaultfontfeatures[\rmfamily]{Ligatures=TeX,Scale=1}





\title{Ejecución condicional en Python}


\shorttitle{CONDICIONALES PYTHON}


\usepackage{etoolbox}



\ccoppy{\textcopyright~2021}



\author{Edison Achalma}



\affiliation{
{Escuela Profesional de Economía, Universidad Nacional de San Cristóbal
de Huamanga}}




\leftheader{Achalma}

\date{2021-06-14}


\abstract{Este abstract será actualizado una vez que se complete el
contenido final del artículo. }

\keywords{keyword1, keyword2}

\authornote{\par{\addORCIDlink{Edison Achalma}{0000-0001-6996-3364}} 
\par{ }
\par{   El autor no tiene conflictos de interés que revelar.    Los
roles de autor se clasificaron utilizando la taxonomía de roles de
colaborador (CRediT; https://credit.niso.org/) de la siguiente
manera:  Edison Achalma:   conceptualización, metodología, análisis
formal, investigación, recursos, curación de
datos, redacción, visualización, supervisión, administración del
proyecto}
\par{La correspondencia relativa a este artículo debe dirigirse a Edison
Achalma, Escuela Profesional de Economía, Universidad Nacional de San
Cristóbal de Huamanga, Portal Independencia N°
57, Ayacucho, AYA 5001, Perú, Email: \href{mailto:elmer.achalma.09@unsch.edu.pe}{elmer.achalma.09@unsch.edu.pe}}
}

\makeatletter
\let\endoldlt\endlongtable
\def\endlongtable{
\hline
\endoldlt
}
\makeatother

\urlstyle{same}



\makeatletter
\@ifpackageloaded{caption}{}{\usepackage{caption}}
\AtBeginDocument{%
\ifdefined\contentsname
  \renewcommand*\contentsname{Tabla de contenidos}
\else
  \newcommand\contentsname{Tabla de contenidos}
\fi
\ifdefined\listfigurename
  \renewcommand*\listfigurename{Lista de Figuras}
\else
  \newcommand\listfigurename{Lista de Figuras}
\fi
\ifdefined\listtablename
  \renewcommand*\listtablename{Lista de Tablas}
\else
  \newcommand\listtablename{Lista de Tablas}
\fi
\ifdefined\figurename
  \renewcommand*\figurename{Figura}
\else
  \newcommand\figurename{Figura}
\fi
\ifdefined\tablename
  \renewcommand*\tablename{Tabla}
\else
  \newcommand\tablename{Tabla}
\fi
}
\@ifpackageloaded{float}{}{\usepackage{float}}
\floatstyle{ruled}
\@ifundefined{c@chapter}{\newfloat{codelisting}{h}{lop}}{\newfloat{codelisting}{h}{lop}[chapter]}
\floatname{codelisting}{Listado}
\newcommand*\listoflistings{\listof{codelisting}{Lista de Listados}}
\makeatother
\makeatletter
\makeatother
\makeatletter
\@ifpackageloaded{caption}{}{\usepackage{caption}}
\@ifpackageloaded{subcaption}{}{\usepackage{subcaption}}
\makeatother
\makeatletter
\@ifpackageloaded{fontawesome5}{}{\usepackage{fontawesome5}}
\makeatother

% From https://tex.stackexchange.com/a/645996/211326
%%% apa7 doesn't want to add appendix section titles in the toc
%%% let's make it do it
\makeatletter
\xpatchcmd{\appendix}
  {\par}
  {\addcontentsline{toc}{section}{\@currentlabelname}\par}
  {}{}
\makeatother

%% Disable longtable counter
%% https://tex.stackexchange.com/a/248395/211326

\usepackage{etoolbox}

\makeatletter
\patchcmd{\LT@caption}
  {\bgroup}
  {\bgroup\global\LTpatch@captiontrue}
  {}{}
\patchcmd{\longtable}
  {\par}
  {\par\global\LTpatch@captionfalse}
  {}{}
\apptocmd{\endlongtable}
  {\ifLTpatch@caption\else\addtocounter{table}{-1}\fi}
  {}{}
\newif\ifLTpatch@caption
\makeatother

\begin{document}

\maketitle


\hypertarget{toc}{}
\tableofcontents
\newpage
\section[Introduction]{Ejecución condicional en Python}

\setcounter{secnumdepth}{3}

\setlength\LTleft{0pt}




En esta cuarta guía exploraremos cómo los programas pueden tomar
decisiones basadas en condiciones. La ejecución condicional es
fundamental para crear análisis económicos automatizados que respondan a
diferentes escenarios y situaciones.

\section{Expresiones booleanas}\label{expresiones-booleanas}

Las expresiones booleanas son expresiones que evalúan a uno de dos
valores posibles: True (verdadero) o False (falso). Son la base de la
toma de decisiones en programación.

\subsection{Valores booleanos}\label{valores-booleanos}

\begin{Shaded}
\begin{Highlighting}[]
\CommentTok{\# Valores booleanos básicos}
\NormalTok{valor\_verdadero }\OperatorTok{=} \VariableTok{True}
\NormalTok{valor\_falso }\OperatorTok{=} \VariableTok{False}

\BuiltInTok{print}\NormalTok{(valor\_verdadero)  }\CommentTok{\# True}
\BuiltInTok{print}\NormalTok{(valor\_falso)      }\CommentTok{\# False}
\BuiltInTok{print}\NormalTok{(}\BuiltInTok{type}\NormalTok{(valor\_verdadero))  }\CommentTok{\# \textless{}class \textquotesingle{}bool\textquotesingle{}\textgreater{}}

\CommentTok{\# En economía}
\NormalTok{economia\_en\_crecimiento }\OperatorTok{=} \VariableTok{True}
\NormalTok{hay\_recesion }\OperatorTok{=} \VariableTok{False}
\NormalTok{meta\_inflacion\_cumplida }\OperatorTok{=} \VariableTok{True}
\end{Highlighting}
\end{Shaded}

\subsection{Función bool()}\label{funciuxf3n-bool}

La función \texttt{bool()} convierte valores a booleanos:

\begin{Shaded}
\begin{Highlighting}[]
\CommentTok{\# Valores que evalúan a False}
\BuiltInTok{print}\NormalTok{(}\BuiltInTok{bool}\NormalTok{(}\DecValTok{0}\NormalTok{))        }\CommentTok{\# False}
\BuiltInTok{print}\NormalTok{(}\BuiltInTok{bool}\NormalTok{(}\FloatTok{0.0}\NormalTok{))      }\CommentTok{\# False}
\BuiltInTok{print}\NormalTok{(}\BuiltInTok{bool}\NormalTok{(}\StringTok{""}\NormalTok{))       }\CommentTok{\# False (string vacío)}
\BuiltInTok{print}\NormalTok{(}\BuiltInTok{bool}\NormalTok{([]))       }\CommentTok{\# False (lista vacía)}
\BuiltInTok{print}\NormalTok{(}\BuiltInTok{bool}\NormalTok{(\{\}))       }\CommentTok{\# False (diccionario vacío)}
\BuiltInTok{print}\NormalTok{(}\BuiltInTok{bool}\NormalTok{(}\VariableTok{None}\NormalTok{))     }\CommentTok{\# False}
\BuiltInTok{print}\NormalTok{(}\BuiltInTok{bool}\NormalTok{(}\VariableTok{False}\NormalTok{))    }\CommentTok{\# False}

\CommentTok{\# Valores que evalúan a True}
\BuiltInTok{print}\NormalTok{(}\BuiltInTok{bool}\NormalTok{(}\DecValTok{1}\NormalTok{))        }\CommentTok{\# True}
\BuiltInTok{print}\NormalTok{(}\BuiltInTok{bool}\NormalTok{(}\OperatorTok{{-}}\DecValTok{1}\NormalTok{))       }\CommentTok{\# True}
\BuiltInTok{print}\NormalTok{(}\BuiltInTok{bool}\NormalTok{(}\FloatTok{0.1}\NormalTok{))      }\CommentTok{\# True}
\BuiltInTok{print}\NormalTok{(}\BuiltInTok{bool}\NormalTok{(}\StringTok{"texto"}\NormalTok{))  }\CommentTok{\# True}
\BuiltInTok{print}\NormalTok{(}\BuiltInTok{bool}\NormalTok{([}\DecValTok{1}\NormalTok{, }\DecValTok{2}\NormalTok{]))   }\CommentTok{\# True}
\BuiltInTok{print}\NormalTok{(}\BuiltInTok{bool}\NormalTok{(\{}\StringTok{\textquotesingle{}a\textquotesingle{}}\NormalTok{: }\DecValTok{1}\NormalTok{\})) }\CommentTok{\# True}
\end{Highlighting}
\end{Shaded}

\subsection{Aplicación económica}\label{aplicaciuxf3n-econuxf3mica}

\begin{Shaded}
\begin{Highlighting}[]
\CommentTok{\# Evaluar datos económicos}
\NormalTok{pib }\OperatorTok{=} \DecValTok{242632}
\NormalTok{inflacion }\OperatorTok{=} \FloatTok{3.5}
\NormalTok{desempleo }\OperatorTok{=} \FloatTok{7.2}

\CommentTok{\# Convertir a booleanos}
\NormalTok{hay\_datos\_pib }\OperatorTok{=} \BuiltInTok{bool}\NormalTok{(pib)}
\NormalTok{hay\_datos\_inflacion }\OperatorTok{=} \BuiltInTok{bool}\NormalTok{(inflacion)}

\BuiltInTok{print}\NormalTok{(}\SpecialStringTok{f"¿Hay datos de PIB? }\SpecialCharTok{\{}\NormalTok{hay\_datos\_pib}\SpecialCharTok{\}}\SpecialStringTok{"}\NormalTok{)}
\BuiltInTok{print}\NormalTok{(}\SpecialStringTok{f"¿Hay datos de inflación? }\SpecialCharTok{\{}\NormalTok{hay\_datos\_inflacion}\SpecialCharTok{\}}\SpecialStringTok{"}\NormalTok{)}

\CommentTok{\# Listas vacías como indicador}
\NormalTok{sectores\_analizados }\OperatorTok{=}\NormalTok{ []}
\NormalTok{regiones\_visitadas }\OperatorTok{=}\NormalTok{ [}\StringTok{\textquotesingle{}Lima\textquotesingle{}}\NormalTok{, }\StringTok{\textquotesingle{}Arequipa\textquotesingle{}}\NormalTok{]}

\BuiltInTok{print}\NormalTok{(}\SpecialStringTok{f"¿Se analizaron sectores? }\SpecialCharTok{\{}\BuiltInTok{bool}\NormalTok{(sectores\_analizados)}\SpecialCharTok{\}}\SpecialStringTok{"}\NormalTok{)  }\CommentTok{\# False}
\BuiltInTok{print}\NormalTok{(}\SpecialStringTok{f"¿Se visitaron regiones? }\SpecialCharTok{\{}\BuiltInTok{bool}\NormalTok{(regiones\_visitadas)}\SpecialCharTok{\}}\SpecialStringTok{"}\NormalTok{)    }\CommentTok{\# True}
\end{Highlighting}
\end{Shaded}

\section{Operadores de comparación}\label{operadores-de-comparaciuxf3n}

Los operadores de comparación comparan dos valores y retornan un
booleano:

\subsection{Operadores básicos}\label{operadores-buxe1sicos}

\begin{Shaded}
\begin{Highlighting}[]
\CommentTok{\# Datos económicos para comparar}
\NormalTok{pib\_2022 }\OperatorTok{=} \DecValTok{240000}
\NormalTok{pib\_2023 }\OperatorTok{=} \DecValTok{242632}
\NormalTok{inflacion\_actual }\OperatorTok{=} \FloatTok{3.5}
\NormalTok{meta\_inflacion }\OperatorTok{=} \FloatTok{3.0}

\CommentTok{\# Igual a (==)}
\BuiltInTok{print}\NormalTok{(pib\_2022 }\OperatorTok{==}\NormalTok{ pib\_2023)  }\CommentTok{\# False}
\BuiltInTok{print}\NormalTok{(}\DecValTok{5} \OperatorTok{==} \FloatTok{5.0}\NormalTok{)              }\CommentTok{\# True (compara valor, no tipo)}

\CommentTok{\# Diferente de (!=)}
\BuiltInTok{print}\NormalTok{(pib\_2022 }\OperatorTok{!=}\NormalTok{ pib\_2023)  }\CommentTok{\# True}
\BuiltInTok{print}\NormalTok{(inflacion\_actual }\OperatorTok{!=}\NormalTok{ meta\_inflacion)  }\CommentTok{\# True}

\CommentTok{\# Mayor que (\textgreater{})}
\BuiltInTok{print}\NormalTok{(pib\_2023 }\OperatorTok{\textgreater{}}\NormalTok{ pib\_2022)   }\CommentTok{\# True}
\BuiltInTok{print}\NormalTok{(inflacion\_actual }\OperatorTok{\textgreater{}}\NormalTok{ meta\_inflacion)  }\CommentTok{\# True}

\CommentTok{\# Menor que (\textless{})}
\BuiltInTok{print}\NormalTok{(pib\_2022 }\OperatorTok{\textless{}}\NormalTok{ pib\_2023)   }\CommentTok{\# True}
\BuiltInTok{print}\NormalTok{(meta\_inflacion }\OperatorTok{\textless{}}\NormalTok{ inflacion\_actual)  }\CommentTok{\# True}

\CommentTok{\# Mayor o igual que (\textgreater{}=)}
\BuiltInTok{print}\NormalTok{(pib\_2023 }\OperatorTok{\textgreater{}=}\NormalTok{ pib\_2022)  }\CommentTok{\# True}
\BuiltInTok{print}\NormalTok{(pib\_2023 }\OperatorTok{\textgreater{}=} \DecValTok{242632}\NormalTok{)    }\CommentTok{\# True}

\CommentTok{\# Menor o igual que (\textless{}=)}
\BuiltInTok{print}\NormalTok{(meta\_inflacion }\OperatorTok{\textless{}=}\NormalTok{ inflacion\_actual)  }\CommentTok{\# True}
\BuiltInTok{print}\NormalTok{(meta\_inflacion }\OperatorTok{\textless{}=} \FloatTok{3.0}\NormalTok{)  }\CommentTok{\# True}
\end{Highlighting}
\end{Shaded}

\subsection{Comparación de strings}\label{comparaciuxf3n-de-strings}

\begin{Shaded}
\begin{Highlighting}[]
\CommentTok{\# Los strings se comparan lexicográficamente}
\NormalTok{pais\_1 }\OperatorTok{=} \StringTok{"Argentina"}
\NormalTok{pais\_2 }\OperatorTok{=} \StringTok{"Perú"}
\NormalTok{pais\_3 }\OperatorTok{=} \StringTok{"Perú"}

\BuiltInTok{print}\NormalTok{(pais\_1 }\OperatorTok{==}\NormalTok{ pais\_2)  }\CommentTok{\# False}
\BuiltInTok{print}\NormalTok{(pais\_2 }\OperatorTok{==}\NormalTok{ pais\_3)  }\CommentTok{\# True}
\BuiltInTok{print}\NormalTok{(pais\_1 }\OperatorTok{\textless{}}\NormalTok{ pais\_2)   }\CommentTok{\# True (orden alfabético)}

\CommentTok{\# Sensibilidad a mayúsculas}
\NormalTok{sector\_1 }\OperatorTok{=} \StringTok{"Minería"}
\NormalTok{sector\_2 }\OperatorTok{=} \StringTok{"minería"}

\BuiltInTok{print}\NormalTok{(sector\_1 }\OperatorTok{==}\NormalTok{ sector\_2)  }\CommentTok{\# False}
\BuiltInTok{print}\NormalTok{(sector\_1.lower() }\OperatorTok{==}\NormalTok{ sector\_2.lower())  }\CommentTok{\# True}
\end{Highlighting}
\end{Shaded}

\subsection{Comparación de listas}\label{comparaciuxf3n-de-listas}

\begin{Shaded}
\begin{Highlighting}[]
\CommentTok{\# Listas se comparan elemento por elemento}
\NormalTok{precios\_2022 }\OperatorTok{=}\NormalTok{ [}\DecValTok{100}\NormalTok{, }\DecValTok{105}\NormalTok{, }\DecValTok{110}\NormalTok{]}
\NormalTok{precios\_2023 }\OperatorTok{=}\NormalTok{ [}\DecValTok{100}\NormalTok{, }\DecValTok{105}\NormalTok{, }\DecValTok{110}\NormalTok{]}
\NormalTok{precios\_2024 }\OperatorTok{=}\NormalTok{ [}\DecValTok{100}\NormalTok{, }\DecValTok{105}\NormalTok{, }\DecValTok{115}\NormalTok{]}

\BuiltInTok{print}\NormalTok{(precios\_2022 }\OperatorTok{==}\NormalTok{ precios\_2023)  }\CommentTok{\# True}
\BuiltInTok{print}\NormalTok{(precios\_2022 }\OperatorTok{==}\NormalTok{ precios\_2024)  }\CommentTok{\# False}
\BuiltInTok{print}\NormalTok{(precios\_2022 }\OperatorTok{\textless{}}\NormalTok{ precios\_2024)   }\CommentTok{\# True (compara primer elemento diferente)}

\CommentTok{\# Orden importa}
\NormalTok{lista\_a }\OperatorTok{=}\NormalTok{ [}\DecValTok{1}\NormalTok{, }\DecValTok{2}\NormalTok{, }\DecValTok{3}\NormalTok{]}
\NormalTok{lista\_b }\OperatorTok{=}\NormalTok{ [}\DecValTok{3}\NormalTok{, }\DecValTok{2}\NormalTok{, }\DecValTok{1}\NormalTok{]}
\BuiltInTok{print}\NormalTok{(lista\_a }\OperatorTok{==}\NormalTok{ lista\_b)  }\CommentTok{\# False}
\end{Highlighting}
\end{Shaded}

\subsection{Ejemplo aplicado: Clasificación de
economías}\label{ejemplo-aplicado-clasificaciuxf3n-de-economuxedas}

\begin{Shaded}
\begin{Highlighting}[]
\CommentTok{\# Clasificar países por PIB per cápita}
\NormalTok{pib\_pc\_peru }\OperatorTok{=} \DecValTok{7196}
\NormalTok{pib\_pc\_chile }\OperatorTok{=} \DecValTok{15400}
\NormalTok{pib\_pc\_argentina }\OperatorTok{=} \DecValTok{10600}

\NormalTok{umbral\_pais\_desarrollado }\OperatorTok{=} \DecValTok{12000}

\BuiltInTok{print}\NormalTok{(}\SpecialStringTok{f"¿Perú es país desarrollado? }\SpecialCharTok{\{}\NormalTok{pib\_pc\_peru }\OperatorTok{\textgreater{}}\NormalTok{ umbral\_pais\_desarrollado}\SpecialCharTok{\}}\SpecialStringTok{"}\NormalTok{)}
\BuiltInTok{print}\NormalTok{(}\SpecialStringTok{f"¿Chile es país desarrollado? }\SpecialCharTok{\{}\NormalTok{pib\_pc\_chile }\OperatorTok{\textgreater{}}\NormalTok{ umbral\_pais\_desarrollado}\SpecialCharTok{\}}\SpecialStringTok{"}\NormalTok{)}
\BuiltInTok{print}\NormalTok{(}\SpecialStringTok{f"¿Argentina es país desarrollado? }\SpecialCharTok{\{}\NormalTok{pib\_pc\_argentina }\OperatorTok{\textgreater{}}\NormalTok{ umbral\_pais\_desarrollado}\SpecialCharTok{\}}\SpecialStringTok{"}\NormalTok{)}

\CommentTok{\# Comparar entre países}
\BuiltInTok{print}\NormalTok{(}\SpecialStringTok{f"}\CharTok{\textbackslash{}n}\SpecialStringTok{¿Chile tiene mayor PIB per cápita que Perú? }\SpecialCharTok{\{}\NormalTok{pib\_pc\_chile }\OperatorTok{\textgreater{}}\NormalTok{ pib\_pc\_peru}\SpecialCharTok{\}}\SpecialStringTok{"}\NormalTok{)}
\BuiltInTok{print}\NormalTok{(}\SpecialStringTok{f"¿Perú y Argentina tienen el mismo PIB per cápita? }\SpecialCharTok{\{}\NormalTok{pib\_pc\_peru }\OperatorTok{==}\NormalTok{ pib\_pc\_argentina}\SpecialCharTok{\}}\SpecialStringTok{"}\NormalTok{)}
\end{Highlighting}
\end{Shaded}

\section{Operadores lógicos}\label{operadores-luxf3gicos}

Los operadores lógicos combinan múltiples expresiones booleanas:

\subsection{Operador AND}\label{operador-and}

El operador \texttt{and} retorna True solo si ambas condiciones son
verdaderas:

\begin{Shaded}
\begin{Highlighting}[]
\CommentTok{\# Tabla de verdad de AND}
\BuiltInTok{print}\NormalTok{(}\VariableTok{True} \KeywordTok{and} \VariableTok{True}\NormalTok{)    }\CommentTok{\# True}
\BuiltInTok{print}\NormalTok{(}\VariableTok{True} \KeywordTok{and} \VariableTok{False}\NormalTok{)   }\CommentTok{\# False}
\BuiltInTok{print}\NormalTok{(}\VariableTok{False} \KeywordTok{and} \VariableTok{True}\NormalTok{)   }\CommentTok{\# False}
\BuiltInTok{print}\NormalTok{(}\VariableTok{False} \KeywordTok{and} \VariableTok{False}\NormalTok{)  }\CommentTok{\# False}

\CommentTok{\# Aplicación económica}
\NormalTok{inflacion }\OperatorTok{=} \FloatTok{3.2}
\NormalTok{desempleo }\OperatorTok{=} \FloatTok{6.8}
\NormalTok{crecimiento }\OperatorTok{=} \FloatTok{2.5}

\CommentTok{\# Economía estable: inflación baja Y desempleo bajo Y crecimiento positivo}
\NormalTok{inflacion\_controlada }\OperatorTok{=}\NormalTok{ inflacion }\OperatorTok{\textless{}} \FloatTok{4.0}
\NormalTok{desempleo\_bajo }\OperatorTok{=}\NormalTok{ desempleo }\OperatorTok{\textless{}} \FloatTok{8.0}
\NormalTok{crecimiento\_positivo }\OperatorTok{=}\NormalTok{ crecimiento }\OperatorTok{\textgreater{}} \DecValTok{0}

\NormalTok{economia\_estable }\OperatorTok{=}\NormalTok{ inflacion\_controlada }\KeywordTok{and}\NormalTok{ desempleo\_bajo }\KeywordTok{and}\NormalTok{ crecimiento\_positivo}
\BuiltInTok{print}\NormalTok{(}\SpecialStringTok{f"¿Economía estable? }\SpecialCharTok{\{}\NormalTok{economia\_estable}\SpecialCharTok{\}}\SpecialStringTok{"}\NormalTok{)  }\CommentTok{\# True}

\CommentTok{\# Ejemplo: Aprobar un crédito}
\NormalTok{ingreso\_mensual }\OperatorTok{=} \DecValTok{5000}
\NormalTok{antiguedad\_laboral }\OperatorTok{=} \DecValTok{3}  \CommentTok{\# años}
\NormalTok{historial\_crediticio }\OperatorTok{=} \StringTok{"bueno"}

\NormalTok{ingreso\_suficiente }\OperatorTok{=}\NormalTok{ ingreso\_mensual }\OperatorTok{\textgreater{}=} \DecValTok{3000}
\NormalTok{estabilidad\_laboral }\OperatorTok{=}\NormalTok{ antiguedad\_laboral }\OperatorTok{\textgreater{}=} \DecValTok{2}
\NormalTok{buen\_historial }\OperatorTok{=}\NormalTok{ historial\_crediticio }\OperatorTok{==} \StringTok{"bueno"}

\NormalTok{aprobar\_credito }\OperatorTok{=}\NormalTok{ ingreso\_suficiente }\KeywordTok{and}\NormalTok{ estabilidad\_laboral }\KeywordTok{and}\NormalTok{ buen\_historial}
\BuiltInTok{print}\NormalTok{(}\SpecialStringTok{f"¿Aprobar crédito? }\SpecialCharTok{\{}\NormalTok{aprobar\_credito}\SpecialCharTok{\}}\SpecialStringTok{"}\NormalTok{)  }\CommentTok{\# True}
\end{Highlighting}
\end{Shaded}

\subsection{Operador OR}\label{operador-or}

El operador \texttt{or} retorna True si al menos una condición es
verdadera:

\begin{Shaded}
\begin{Highlighting}[]
\CommentTok{\# Tabla de verdad de OR}
\BuiltInTok{print}\NormalTok{(}\VariableTok{True} \KeywordTok{or} \VariableTok{True}\NormalTok{)     }\CommentTok{\# True}
\BuiltInTok{print}\NormalTok{(}\VariableTok{True} \KeywordTok{or} \VariableTok{False}\NormalTok{)    }\CommentTok{\# True}
\BuiltInTok{print}\NormalTok{(}\VariableTok{False} \KeywordTok{or} \VariableTok{True}\NormalTok{)    }\CommentTok{\# True}
\BuiltInTok{print}\NormalTok{(}\VariableTok{False} \KeywordTok{or} \VariableTok{False}\NormalTok{)   }\CommentTok{\# False}

\CommentTok{\# Aplicación económica: Señales de alerta}
\NormalTok{inflacion }\OperatorTok{=} \FloatTok{5.5}
\NormalTok{desempleo }\OperatorTok{=} \FloatTok{9.2}
\NormalTok{deficit\_fiscal }\OperatorTok{=} \FloatTok{4.8}

\NormalTok{inflacion\_alta }\OperatorTok{=}\NormalTok{ inflacion }\OperatorTok{\textgreater{}} \FloatTok{5.0}
\NormalTok{desempleo\_alto }\OperatorTok{=}\NormalTok{ desempleo }\OperatorTok{\textgreater{}} \FloatTok{8.0}
\NormalTok{deficit\_alto }\OperatorTok{=}\NormalTok{ deficit\_fiscal }\OperatorTok{\textgreater{}} \FloatTok{4.0}

\NormalTok{alerta\_economica }\OperatorTok{=}\NormalTok{ inflacion\_alta }\KeywordTok{or}\NormalTok{ desempleo\_alto }\KeywordTok{or}\NormalTok{ deficit\_alto}
\BuiltInTok{print}\NormalTok{(}\SpecialStringTok{f"¿Hay alerta económica? }\SpecialCharTok{\{}\NormalTok{alerta\_economica}\SpecialCharTok{\}}\SpecialStringTok{"}\NormalTok{)  }\CommentTok{\# True}

\CommentTok{\# Ejemplo: Elegibilidad para programa social}
\NormalTok{ingreso\_familiar }\OperatorTok{=} \DecValTok{1200}
\NormalTok{numero\_dependientes }\OperatorTok{=} \DecValTok{3}
\NormalTok{zona\_rural }\OperatorTok{=} \VariableTok{True}

\NormalTok{ingreso\_bajo }\OperatorTok{=}\NormalTok{ ingreso\_familiar }\OperatorTok{\textless{}} \DecValTok{1500}
\NormalTok{familia\_numerosa }\OperatorTok{=}\NormalTok{ numero\_dependientes }\OperatorTok{\textgreater{}=} \DecValTok{3}

\NormalTok{elegible }\OperatorTok{=}\NormalTok{ ingreso\_bajo }\KeywordTok{or}\NormalTok{ familia\_numerosa }\KeywordTok{or}\NormalTok{ zona\_rural}
\BuiltInTok{print}\NormalTok{(}\SpecialStringTok{f"¿Elegible para programa? }\SpecialCharTok{\{}\NormalTok{elegible}\SpecialCharTok{\}}\SpecialStringTok{"}\NormalTok{)  }\CommentTok{\# True}
\end{Highlighting}
\end{Shaded}

\subsection{Operador NOT}\label{operador-not}

El operador \texttt{not} invierte el valor booleano:

\begin{Shaded}
\begin{Highlighting}[]
\CommentTok{\# Tabla de verdad de NOT}
\BuiltInTok{print}\NormalTok{(}\KeywordTok{not} \VariableTok{True}\NormalTok{)   }\CommentTok{\# False}
\BuiltInTok{print}\NormalTok{(}\KeywordTok{not} \VariableTok{False}\NormalTok{)  }\CommentTok{\# True}

\CommentTok{\# Aplicación económica}
\NormalTok{en\_recesion }\OperatorTok{=} \VariableTok{False}
\NormalTok{tiene\_empleo }\OperatorTok{=} \VariableTok{True}
\NormalTok{deuda\_impagable }\OperatorTok{=} \VariableTok{False}

\BuiltInTok{print}\NormalTok{(}\SpecialStringTok{f"¿No está en recesión? }\SpecialCharTok{\{}\KeywordTok{not}\NormalTok{ en\_recesion}\SpecialCharTok{\}}\SpecialStringTok{"}\NormalTok{)  }\CommentTok{\# True}
\BuiltInTok{print}\NormalTok{(}\SpecialStringTok{f"¿Está desempleado? }\SpecialCharTok{\{}\KeywordTok{not}\NormalTok{ tiene\_empleo}\SpecialCharTok{\}}\SpecialStringTok{"}\NormalTok{)    }\CommentTok{\# False}
\BuiltInTok{print}\NormalTok{(}\SpecialStringTok{f"¿Deuda manejable? }\SpecialCharTok{\{}\KeywordTok{not}\NormalTok{ deuda\_impagable}\SpecialCharTok{\}}\SpecialStringTok{"}\NormalTok{)  }\CommentTok{\# True}

\CommentTok{\# Ejemplo: Verificar condiciones negativas}
\NormalTok{cumple\_meta\_inflacion }\OperatorTok{=} \VariableTok{False}
\NormalTok{supera\_crecimiento\_esperado }\OperatorTok{=} \VariableTok{True}

\NormalTok{no\_cumple\_meta }\OperatorTok{=} \KeywordTok{not}\NormalTok{ cumple\_meta\_inflacion}
\BuiltInTok{print}\NormalTok{(}\SpecialStringTok{f"¿Necesita ajuste de política? }\SpecialCharTok{\{}\NormalTok{no\_cumple\_meta}\SpecialCharTok{\}}\SpecialStringTok{"}\NormalTok{)  }\CommentTok{\# True}
\end{Highlighting}
\end{Shaded}

\subsection{Combinación de
operadores}\label{combinaciuxf3n-de-operadores}

\begin{Shaded}
\begin{Highlighting}[]
\CommentTok{\# Los operadores se pueden combinar}
\CommentTok{\# Precedencia: not \textgreater{} and \textgreater{} or}

\CommentTok{\# Análisis de inversión}
\NormalTok{rentabilidad }\OperatorTok{=} \FloatTok{8.5}
\NormalTok{riesgo }\OperatorTok{=} \StringTok{"medio"}
\NormalTok{liquidez }\OperatorTok{=} \StringTok{"alta"}
\NormalTok{plazo\_anos }\OperatorTok{=} \DecValTok{3}

\CommentTok{\# Condiciones}
\NormalTok{rentabilidad\_aceptable }\OperatorTok{=}\NormalTok{ rentabilidad }\OperatorTok{\textgreater{}} \FloatTok{6.0}
\NormalTok{riesgo\_aceptable }\OperatorTok{=}\NormalTok{ riesgo }\KeywordTok{in}\NormalTok{ [}\StringTok{"bajo"}\NormalTok{, }\StringTok{"medio"}\NormalTok{]}
\NormalTok{liquidez\_aceptable }\OperatorTok{=}\NormalTok{ liquidez }\KeywordTok{in}\NormalTok{ [}\StringTok{"media"}\NormalTok{, }\StringTok{"alta"}\NormalTok{]}
\NormalTok{plazo\_adecuado }\OperatorTok{=}\NormalTok{ plazo\_anos }\OperatorTok{\textless{}=} \DecValTok{5}

\CommentTok{\# Decisión de inversión}
\NormalTok{invertir }\OperatorTok{=}\NormalTok{ (rentabilidad\_aceptable }\KeywordTok{and}\NormalTok{ riesgo\_aceptable) }\KeywordTok{and} \OperatorTok{\textbackslash{}}
\NormalTok{           (liquidez\_aceptable }\KeywordTok{or}\NormalTok{ plazo\_adecuado)}

\BuiltInTok{print}\NormalTok{(}\SpecialStringTok{f"¿Realizar inversión? }\SpecialCharTok{\{}\NormalTok{invertir}\SpecialCharTok{\}}\SpecialStringTok{"}\NormalTok{)}

\CommentTok{\# Uso de paréntesis para claridad}
\NormalTok{inflacion }\OperatorTok{=} \FloatTok{4.2}
\NormalTok{desempleo }\OperatorTok{=} \FloatTok{7.5}
\NormalTok{reservas }\OperatorTok{=} \DecValTok{75000}

\NormalTok{situacion\_critica }\OperatorTok{=}\NormalTok{ (inflacion }\OperatorTok{\textgreater{}} \FloatTok{5.0} \KeywordTok{and}\NormalTok{ desempleo }\OperatorTok{\textgreater{}} \FloatTok{10.0}\NormalTok{) }\KeywordTok{or} \OperatorTok{\textbackslash{}}
\NormalTok{                    (reservas }\OperatorTok{\textless{}} \DecValTok{50000}\NormalTok{)}

\BuiltInTok{print}\NormalTok{(}\SpecialStringTok{f"¿Situación crítica? }\SpecialCharTok{\{}\NormalTok{situacion\_critica}\SpecialCharTok{\}}\SpecialStringTok{"}\NormalTok{)  }\CommentTok{\# False}
\end{Highlighting}
\end{Shaded}

\section{Ejecución condicional
simple}\label{ejecuciuxf3n-condicional-simple}

La estructura \texttt{if} ejecuta código solo si una condición es
verdadera:

\subsection{Sintaxis básica}\label{sintaxis-buxe1sica}

\begin{Shaded}
\begin{Highlighting}[]
\CommentTok{\# Estructura básica}
\NormalTok{condicion }\OperatorTok{=} \VariableTok{True}

\ControlFlowTok{if}\NormalTok{ condicion:}
    \BuiltInTok{print}\NormalTok{(}\StringTok{"La condición es verdadera"}\NormalTok{)}
    \BuiltInTok{print}\NormalTok{(}\StringTok{"Este código se ejecuta"}\NormalTok{)}

\CommentTok{\# Sin indentación, no es parte del if}
\BuiltInTok{print}\NormalTok{(}\StringTok{"Este código siempre se ejecuta"}\NormalTok{)}
\end{Highlighting}
\end{Shaded}

\subsection{Ejemplos económicos}\label{ejemplos-econuxf3micos}

\begin{Shaded}
\begin{Highlighting}[]
\CommentTok{\# Ejemplo 1: Alerta de inflación}
\NormalTok{inflacion\_mensual }\OperatorTok{=} \FloatTok{0.6}

\ControlFlowTok{if}\NormalTok{ inflacion\_mensual }\OperatorTok{\textgreater{}} \FloatTok{0.5}\NormalTok{:}
    \BuiltInTok{print}\NormalTok{(}\StringTok{"ALERTA: Inflación mensual supera el 0.5\%"}\NormalTok{)}
    \BuiltInTok{print}\NormalTok{(}\SpecialStringTok{f"Inflación registrada: }\SpecialCharTok{\{}\NormalTok{inflacion\_mensual}\SpecialCharTok{\}}\SpecialStringTok{\%"}\NormalTok{)}
    \BuiltInTok{print}\NormalTok{(}\StringTok{"Se recomienda revisar política monetaria"}\NormalTok{)}

\CommentTok{\# Ejemplo 2: Verificar meta de crecimiento}
\NormalTok{crecimiento\_pib }\OperatorTok{=} \FloatTok{3.2}
\NormalTok{meta\_crecimiento }\OperatorTok{=} \FloatTok{3.0}

\ControlFlowTok{if}\NormalTok{ crecimiento\_pib }\OperatorTok{\textgreater{}=}\NormalTok{ meta\_crecimiento:}
    \BuiltInTok{print}\NormalTok{(}\StringTok{"Meta de crecimiento alcanzada"}\NormalTok{)}
    \BuiltInTok{print}\NormalTok{(}\SpecialStringTok{f"Crecimiento: }\SpecialCharTok{\{}\NormalTok{crecimiento\_pib}\SpecialCharTok{\}}\SpecialStringTok{\%"}\NormalTok{)}
    \BuiltInTok{print}\NormalTok{(}\SpecialStringTok{f"Meta: }\SpecialCharTok{\{}\NormalTok{meta\_crecimiento}\SpecialCharTok{\}}\SpecialStringTok{\%"}\NormalTok{)}

\CommentTok{\# Ejemplo 3: Control de inventario}
\NormalTok{stock\_actual }\OperatorTok{=} \DecValTok{8}
\NormalTok{stock\_minimo }\OperatorTok{=} \DecValTok{10}

\ControlFlowTok{if}\NormalTok{ stock\_actual }\OperatorTok{\textless{}}\NormalTok{ stock\_minimo:}
    \BuiltInTok{print}\NormalTok{(}\StringTok{"ALERTA: Stock bajo el mínimo"}\NormalTok{)}
\NormalTok{    faltante }\OperatorTok{=}\NormalTok{ stock\_minimo }\OperatorTok{{-}}\NormalTok{ stock\_actual}
    \BuiltInTok{print}\NormalTok{(}\SpecialStringTok{f"Se requiere reponer }\SpecialCharTok{\{}\NormalTok{faltante}\SpecialCharTok{\}}\SpecialStringTok{ unidades"}\NormalTok{)}
\end{Highlighting}
\end{Shaded}

\subsection{Importancia de la
indentación}\label{importancia-de-la-indentaciuxf3n}

\begin{Shaded}
\begin{Highlighting}[]
\CommentTok{\# Python usa indentación para delimitar bloques}
\NormalTok{x }\OperatorTok{=} \DecValTok{10}

\ControlFlowTok{if}\NormalTok{ x }\OperatorTok{\textgreater{}} \DecValTok{5}\NormalTok{:}
    \BuiltInTok{print}\NormalTok{(}\StringTok{"x es mayor que 5"}\NormalTok{)}
    \BuiltInTok{print}\NormalTok{(}\StringTok{"Este también es parte del if"}\NormalTok{)}
\BuiltInTok{print}\NormalTok{(}\StringTok{"Este no es parte del if"}\NormalTok{)}

\CommentTok{\# Error de indentación (comentado para evitar error)}
\CommentTok{\# if x \textgreater{} 5:}
\CommentTok{\# print("Error: no hay indentación")}

\CommentTok{\# Indentación inconsistente (comentado)}
\CommentTok{\# if x \textgreater{} 5:}
\CommentTok{\#     print("4 espacios")}
\CommentTok{\#       print("6 espacios {-} Error!")}
\end{Highlighting}
\end{Shaded}

\section{Ejecución alternativa}\label{ejecuciuxf3n-alternativa}

La estructura \texttt{if-else} ejecuta un bloque si la condición es
verdadera, y otro si es falsa:

\subsection{Sintaxis if-else}\label{sintaxis-if-else}

\begin{Shaded}
\begin{Highlighting}[]
\CommentTok{\# Estructura básica}
\NormalTok{condicion }\OperatorTok{=} \VariableTok{True}

\ControlFlowTok{if}\NormalTok{ condicion:}
    \BuiltInTok{print}\NormalTok{(}\StringTok{"Condición verdadera"}\NormalTok{)}
\ControlFlowTok{else}\NormalTok{:}
    \BuiltInTok{print}\NormalTok{(}\StringTok{"Condición falsa"}\NormalTok{)}
\end{Highlighting}
\end{Shaded}

\subsection{Ejemplos económicos}\label{ejemplos-econuxf3micos-1}

\begin{Shaded}
\begin{Highlighting}[]
\CommentTok{\# Ejemplo 1: Clasificar inflación}
\NormalTok{inflacion }\OperatorTok{=} \FloatTok{3.5}
\NormalTok{meta }\OperatorTok{=} \FloatTok{3.0}

\ControlFlowTok{if}\NormalTok{ inflacion }\OperatorTok{\textless{}=}\NormalTok{ meta:}
    \BuiltInTok{print}\NormalTok{(}\StringTok{"Inflación bajo control"}\NormalTok{)}
    \BuiltInTok{print}\NormalTok{(}\SpecialStringTok{f"Inflación: }\SpecialCharTok{\{}\NormalTok{inflacion}\SpecialCharTok{\}}\SpecialStringTok{\%"}\NormalTok{)}
    \BuiltInTok{print}\NormalTok{(}\SpecialStringTok{f"Meta: }\SpecialCharTok{\{}\NormalTok{meta}\SpecialCharTok{\}}\SpecialStringTok{\%"}\NormalTok{)}
\ControlFlowTok{else}\NormalTok{:}
    \BuiltInTok{print}\NormalTok{(}\StringTok{"Inflación por encima de la meta"}\NormalTok{)}
\NormalTok{    exceso }\OperatorTok{=}\NormalTok{ inflacion }\OperatorTok{{-}}\NormalTok{ meta}
    \BuiltInTok{print}\NormalTok{(}\SpecialStringTok{f"Exceso: }\SpecialCharTok{\{}\NormalTok{exceso}\SpecialCharTok{\}}\SpecialStringTok{\%"}\NormalTok{)}
    \BuiltInTok{print}\NormalTok{(}\StringTok{"Se requiere ajuste en política monetaria"}\NormalTok{)}

\CommentTok{\# Ejemplo 2: Determinar superávit o déficit}
\NormalTok{ingresos }\OperatorTok{=} \DecValTok{1500000}
\NormalTok{gastos }\OperatorTok{=} \DecValTok{1450000}

\ControlFlowTok{if}\NormalTok{ ingresos }\OperatorTok{\textgreater{}=}\NormalTok{ gastos:}
\NormalTok{    superavit }\OperatorTok{=}\NormalTok{ ingresos }\OperatorTok{{-}}\NormalTok{ gastos}
    \BuiltInTok{print}\NormalTok{(}\SpecialStringTok{f"Superávit fiscal: S/ }\SpecialCharTok{\{}\NormalTok{superavit}\SpecialCharTok{:,\}}\SpecialStringTok{"}\NormalTok{)}
    \BuiltInTok{print}\NormalTok{(}\StringTok{"Situación fiscal favorable"}\NormalTok{)}
\ControlFlowTok{else}\NormalTok{:}
\NormalTok{    deficit }\OperatorTok{=}\NormalTok{ gastos }\OperatorTok{{-}}\NormalTok{ ingresos}
    \BuiltInTok{print}\NormalTok{(}\SpecialStringTok{f"Déficit fiscal: S/ }\SpecialCharTok{\{}\NormalTok{deficit}\SpecialCharTok{:,\}}\SpecialStringTok{"}\NormalTok{)}
    \BuiltInTok{print}\NormalTok{(}\StringTok{"Se requiere ajuste fiscal"}\NormalTok{)}

\CommentTok{\# Ejemplo 3: Calificación de riesgo crediticio}
\NormalTok{puntaje\_crediticio }\OperatorTok{=} \DecValTok{720}
\NormalTok{umbral }\OperatorTok{=} \DecValTok{700}

\ControlFlowTok{if}\NormalTok{ puntaje\_crediticio }\OperatorTok{\textgreater{}=}\NormalTok{ umbral:}
\NormalTok{    tasa\_interes }\OperatorTok{=} \FloatTok{8.5}
    \BuiltInTok{print}\NormalTok{(}\StringTok{"Crédito aprobado"}\NormalTok{)}
    \BuiltInTok{print}\NormalTok{(}\SpecialStringTok{f"Tasa de interés: }\SpecialCharTok{\{}\NormalTok{tasa\_interes}\SpecialCharTok{\}}\SpecialStringTok{\%"}\NormalTok{)}
\ControlFlowTok{else}\NormalTok{:}
    \BuiltInTok{print}\NormalTok{(}\StringTok{"Crédito denegado"}\NormalTok{)}
    \BuiltInTok{print}\NormalTok{(}\StringTok{"Puntaje insuficiente"}\NormalTok{)}
    \BuiltInTok{print}\NormalTok{(}\SpecialStringTok{f"Puntaje actual: }\SpecialCharTok{\{}\NormalTok{puntaje\_crediticio}\SpecialCharTok{\}}\SpecialStringTok{"}\NormalTok{)}
    \BuiltInTok{print}\NormalTok{(}\SpecialStringTok{f"Puntaje requerido: }\SpecialCharTok{\{}\NormalTok{umbral}\SpecialCharTok{\}}\SpecialStringTok{"}\NormalTok{)}
\end{Highlighting}
\end{Shaded}

\subsection{Operador ternario}\label{operador-ternario}

Python permite escribir if-else en una línea:

\begin{Shaded}
\begin{Highlighting}[]
\CommentTok{\# Sintaxis: valor\_si\_true if condicion else valor\_si\_false}

\CommentTok{\# Ejemplo 1: Clasificar país}
\NormalTok{pib\_per\_capita }\OperatorTok{=} \DecValTok{15000}
\NormalTok{clasificacion }\OperatorTok{=} \StringTok{"Desarrollado"} \ControlFlowTok{if}\NormalTok{ pib\_per\_capita }\OperatorTok{\textgreater{}} \DecValTok{12000} \ControlFlowTok{else} \StringTok{"En desarrollo"}
\BuiltInTok{print}\NormalTok{(}\SpecialStringTok{f"Clasificación: }\SpecialCharTok{\{}\NormalTok{clasificacion}\SpecialCharTok{\}}\SpecialStringTok{"}\NormalTok{)}

\CommentTok{\# Ejemplo 2: Determinar signo}
\NormalTok{saldo }\OperatorTok{=} \DecValTok{5000}
\NormalTok{signo }\OperatorTok{=} \StringTok{"+"} \ControlFlowTok{if}\NormalTok{ saldo }\OperatorTok{\textgreater{}=} \DecValTok{0} \ControlFlowTok{else} \StringTok{"{-}"}
\BuiltInTok{print}\NormalTok{(}\SpecialStringTok{f"Saldo: }\SpecialCharTok{\{}\NormalTok{signo}\SpecialCharTok{\}}\SpecialStringTok{S/ }\SpecialCharTok{\{}\BuiltInTok{abs}\NormalTok{(saldo)}\SpecialCharTok{:,\}}\SpecialStringTok{"}\NormalTok{)}

\CommentTok{\# Ejemplo 3: Calcular descuento}
\NormalTok{precio }\OperatorTok{=} \DecValTok{1000}
\NormalTok{cantidad }\OperatorTok{=} \DecValTok{15}
\NormalTok{descuento }\OperatorTok{=} \FloatTok{0.10} \ControlFlowTok{if}\NormalTok{ cantidad }\OperatorTok{\textgreater{}=} \DecValTok{10} \ControlFlowTok{else} \DecValTok{0}
\NormalTok{precio\_final }\OperatorTok{=}\NormalTok{ precio }\OperatorTok{*}\NormalTok{ (}\DecValTok{1} \OperatorTok{{-}}\NormalTok{ descuento)}
\BuiltInTok{print}\NormalTok{(}\SpecialStringTok{f"Precio final: S/ }\SpecialCharTok{\{}\NormalTok{precio\_final}\SpecialCharTok{:.2f\}}\SpecialStringTok{"}\NormalTok{)}
\end{Highlighting}
\end{Shaded}

\section{Condicionales encadenados}\label{condicionales-encadenados}

La estructura \texttt{if-elif-else} permite evaluar múltiples
condiciones en secuencia:

\subsection{Sintaxis if-elif-else}\label{sintaxis-if-elif-else}

\begin{Shaded}
\begin{Highlighting}[]
\CommentTok{\# Estructura básica}
\NormalTok{x }\OperatorTok{=} \DecValTok{10}

\ControlFlowTok{if}\NormalTok{ x }\OperatorTok{\textless{}} \DecValTok{0}\NormalTok{:}
    \BuiltInTok{print}\NormalTok{(}\StringTok{"Negativo"}\NormalTok{)}
\ControlFlowTok{elif}\NormalTok{ x }\OperatorTok{==} \DecValTok{0}\NormalTok{:}
    \BuiltInTok{print}\NormalTok{(}\StringTok{"Cero"}\NormalTok{)}
\ControlFlowTok{elif}\NormalTok{ x }\OperatorTok{\textless{}} \DecValTok{10}\NormalTok{:}
    \BuiltInTok{print}\NormalTok{(}\StringTok{"Positivo menor que 10"}\NormalTok{)}
\ControlFlowTok{else}\NormalTok{:}
    \BuiltInTok{print}\NormalTok{(}\StringTok{"10 o mayor"}\NormalTok{)}
\end{Highlighting}
\end{Shaded}

\subsection{Ejemplos económicos}\label{ejemplos-econuxf3micos-2}

\begin{Shaded}
\begin{Highlighting}[]
\CommentTok{\# Ejemplo 1: Clasificación de inflación}
\NormalTok{inflacion }\OperatorTok{=} \FloatTok{4.5}

\ControlFlowTok{if}\NormalTok{ inflacion }\OperatorTok{\textless{}} \FloatTok{2.0}\NormalTok{:}
\NormalTok{    clasificacion }\OperatorTok{=} \StringTok{"Muy baja"}
\NormalTok{    comentario }\OperatorTok{=} \StringTok{"Riesgo de deflación"}
\ControlFlowTok{elif}\NormalTok{ inflacion }\OperatorTok{\textless{}} \FloatTok{3.0}\NormalTok{:}
\NormalTok{    clasificacion }\OperatorTok{=} \StringTok{"Baja"}
\NormalTok{    comentario }\OperatorTok{=} \StringTok{"Dentro del objetivo"}
\ControlFlowTok{elif}\NormalTok{ inflacion }\OperatorTok{\textless{}} \FloatTok{4.0}\NormalTok{:}
\NormalTok{    clasificacion }\OperatorTok{=} \StringTok{"Moderada"}
\NormalTok{    comentario }\OperatorTok{=} \StringTok{"Ligeramente por encima del objetivo"}
\ControlFlowTok{elif}\NormalTok{ inflacion }\OperatorTok{\textless{}} \FloatTok{6.0}\NormalTok{:}
\NormalTok{    clasificacion }\OperatorTok{=} \StringTok{"Alta"}
\NormalTok{    comentario }\OperatorTok{=} \StringTok{"Requiere medidas correctivas"}
\ControlFlowTok{else}\NormalTok{:}
\NormalTok{    clasificacion }\OperatorTok{=} \StringTok{"Muy alta"}
\NormalTok{    comentario }\OperatorTok{=} \StringTok{"Situación crítica"}

\BuiltInTok{print}\NormalTok{(}\SpecialStringTok{f"Inflación: }\SpecialCharTok{\{}\NormalTok{inflacion}\SpecialCharTok{\}}\SpecialStringTok{\%"}\NormalTok{)}
\BuiltInTok{print}\NormalTok{(}\SpecialStringTok{f"Clasificación: }\SpecialCharTok{\{}\NormalTok{clasificacion}\SpecialCharTok{\}}\SpecialStringTok{"}\NormalTok{)}
\BuiltInTok{print}\NormalTok{(}\SpecialStringTok{f"Comentario: }\SpecialCharTok{\{}\NormalTok{comentario}\SpecialCharTok{\}}\SpecialStringTok{"}\NormalTok{)}

\CommentTok{\# Ejemplo 2: Escala de riesgo crediticio}
\NormalTok{puntaje }\OperatorTok{=} \DecValTok{680}

\ControlFlowTok{if}\NormalTok{ puntaje }\OperatorTok{\textgreater{}=} \DecValTok{800}\NormalTok{:}
\NormalTok{    riesgo }\OperatorTok{=} \StringTok{"Muy bajo"}
\NormalTok{    tasa }\OperatorTok{=} \FloatTok{7.5}
\ControlFlowTok{elif}\NormalTok{ puntaje }\OperatorTok{\textgreater{}=} \DecValTok{740}\NormalTok{:}
\NormalTok{    riesgo }\OperatorTok{=} \StringTok{"Bajo"}
\NormalTok{    tasa }\OperatorTok{=} \FloatTok{9.0}
\ControlFlowTok{elif}\NormalTok{ puntaje }\OperatorTok{\textgreater{}=} \DecValTok{670}\NormalTok{:}
\NormalTok{    riesgo }\OperatorTok{=} \StringTok{"Medio"}
\NormalTok{    tasa }\OperatorTok{=} \FloatTok{12.0}
\ControlFlowTok{elif}\NormalTok{ puntaje }\OperatorTok{\textgreater{}=} \DecValTok{580}\NormalTok{:}
\NormalTok{    riesgo }\OperatorTok{=} \StringTok{"Alto"}
\NormalTok{    tasa }\OperatorTok{=} \FloatTok{16.0}
\ControlFlowTok{else}\NormalTok{:}
\NormalTok{    riesgo }\OperatorTok{=} \StringTok{"Muy alto"}
\NormalTok{    tasa }\OperatorTok{=} \VariableTok{None}

\BuiltInTok{print}\NormalTok{(}\SpecialStringTok{f"Puntaje crediticio: }\SpecialCharTok{\{}\NormalTok{puntaje}\SpecialCharTok{\}}\SpecialStringTok{"}\NormalTok{)}
\BuiltInTok{print}\NormalTok{(}\SpecialStringTok{f"Nivel de riesgo: }\SpecialCharTok{\{}\NormalTok{riesgo}\SpecialCharTok{\}}\SpecialStringTok{"}\NormalTok{)}
\ControlFlowTok{if}\NormalTok{ tasa:}
    \BuiltInTok{print}\NormalTok{(}\SpecialStringTok{f"Tasa de interés: }\SpecialCharTok{\{}\NormalTok{tasa}\SpecialCharTok{\}}\SpecialStringTok{\%"}\NormalTok{)}
\ControlFlowTok{else}\NormalTok{:}
    \BuiltInTok{print}\NormalTok{(}\StringTok{"Crédito no disponible"}\NormalTok{)}

\CommentTok{\# Ejemplo 3: Clasificación de países por ingreso}
\NormalTok{ingreso\_per\_capita }\OperatorTok{=} \DecValTok{4500}

\ControlFlowTok{if}\NormalTok{ ingreso\_per\_capita }\OperatorTok{\textless{}} \DecValTok{1045}\NormalTok{:}
\NormalTok{    categoria }\OperatorTok{=} \StringTok{"Ingreso bajo"}
\NormalTok{    grupo }\OperatorTok{=} \StringTok{"Países menos desarrollados"}
\ControlFlowTok{elif}\NormalTok{ ingreso\_per\_capita }\OperatorTok{\textless{}} \DecValTok{4095}\NormalTok{:}
\NormalTok{    categoria }\OperatorTok{=} \StringTok{"Ingreso medio{-}bajo"}
\NormalTok{    grupo }\OperatorTok{=} \StringTok{"Economías emergentes"}
\ControlFlowTok{elif}\NormalTok{ ingreso\_per\_capita }\OperatorTok{\textless{}} \DecValTok{12695}\NormalTok{:}
\NormalTok{    categoria }\OperatorTok{=} \StringTok{"Ingreso medio{-}alto"}
\NormalTok{    grupo }\OperatorTok{=} \StringTok{"Economías en transición"}
\ControlFlowTok{else}\NormalTok{:}
\NormalTok{    categoria }\OperatorTok{=} \StringTok{"Ingreso alto"}
\NormalTok{    grupo }\OperatorTok{=} \StringTok{"Países desarrollados"}

\BuiltInTok{print}\NormalTok{(}\SpecialStringTok{f"Ingreso per cápita: USD }\SpecialCharTok{\{}\NormalTok{ingreso\_per\_capita}\SpecialCharTok{:,\}}\SpecialStringTok{"}\NormalTok{)}
\BuiltInTok{print}\NormalTok{(}\SpecialStringTok{f"Categoría: }\SpecialCharTok{\{}\NormalTok{categoria}\SpecialCharTok{\}}\SpecialStringTok{"}\NormalTok{)}
\BuiltInTok{print}\NormalTok{(}\SpecialStringTok{f"Grupo: }\SpecialCharTok{\{}\NormalTok{grupo}\SpecialCharTok{\}}\SpecialStringTok{"}\NormalTok{)}
\end{Highlighting}
\end{Shaded}

\subsection{Orden de evaluación}\label{orden-de-evaluaciuxf3n}

\begin{Shaded}
\begin{Highlighting}[]
\CommentTok{\# El orden importa: se evalúa de arriba hacia abajo}
\CommentTok{\# Solo se ejecuta el primer bloque que cumple la condición}

\NormalTok{nota }\OperatorTok{=} \DecValTok{85}

\ControlFlowTok{if}\NormalTok{ nota }\OperatorTok{\textgreater{}=} \DecValTok{90}\NormalTok{:}
\NormalTok{    calificacion }\OperatorTok{=} \StringTok{"Excelente"}
\ControlFlowTok{elif}\NormalTok{ nota }\OperatorTok{\textgreater{}=} \DecValTok{80}\NormalTok{:}
\NormalTok{    calificacion }\OperatorTok{=} \StringTok{"Muy bueno"}
\ControlFlowTok{elif}\NormalTok{ nota }\OperatorTok{\textgreater{}=} \DecValTok{70}\NormalTok{:}
\NormalTok{    calificacion }\OperatorTok{=} \StringTok{"Bueno"}
\ControlFlowTok{elif}\NormalTok{ nota }\OperatorTok{\textgreater{}=} \DecValTok{60}\NormalTok{:}
\NormalTok{    calificacion }\OperatorTok{=} \StringTok{"Regular"}
\ControlFlowTok{else}\NormalTok{:}
\NormalTok{    calificacion }\OperatorTok{=} \StringTok{"Insuficiente"}

\BuiltInTok{print}\NormalTok{(}\SpecialStringTok{f"Nota: }\SpecialCharTok{\{}\NormalTok{nota}\SpecialCharTok{\}}\SpecialStringTok{"}\NormalTok{)}
\BuiltInTok{print}\NormalTok{(}\SpecialStringTok{f"Calificación: }\SpecialCharTok{\{}\NormalTok{calificacion}\SpecialCharTok{\}}\SpecialStringTok{"}\NormalTok{)}

\CommentTok{\# Si cambiamos el orden, el resultado puede ser diferente}
\CommentTok{\# (Ejemplo de lo que NO hacer)}
\ControlFlowTok{if}\NormalTok{ nota }\OperatorTok{\textgreater{}=} \DecValTok{60}\NormalTok{:  }\CommentTok{\# Este se evalúa primero}
\NormalTok{    calificacion\_incorrecta }\OperatorTok{=} \StringTok{"Regular"}
\ControlFlowTok{elif}\NormalTok{ nota }\OperatorTok{\textgreater{}=} \DecValTok{90}\NormalTok{:  }\CommentTok{\# Nunca se alcanza}
\NormalTok{    calificacion\_incorrecta }\OperatorTok{=} \StringTok{"Excelente"}

\BuiltInTok{print}\NormalTok{(}\SpecialStringTok{f"Calificación incorrecta: }\SpecialCharTok{\{}\NormalTok{calificacion\_incorrecta}\SpecialCharTok{\}}\SpecialStringTok{"}\NormalTok{)  }\CommentTok{\# Regular}
\end{Highlighting}
\end{Shaded}

\section{Condicionales anidados}\label{condicionales-anidados}

Los condicionales pueden contener otros condicionales dentro:

\subsection{Estructura básica}\label{estructura-buxe1sica}

\begin{Shaded}
\begin{Highlighting}[]
\CommentTok{\# Condicional anidado simple}
\NormalTok{x }\OperatorTok{=} \DecValTok{10}

\ControlFlowTok{if}\NormalTok{ x }\OperatorTok{\textgreater{}} \DecValTok{0}\NormalTok{:}
    \BuiltInTok{print}\NormalTok{(}\StringTok{"x es positivo"}\NormalTok{)}
    \ControlFlowTok{if}\NormalTok{ x }\OperatorTok{\textgreater{}} \DecValTok{5}\NormalTok{:}
        \BuiltInTok{print}\NormalTok{(}\StringTok{"x es mayor que 5"}\NormalTok{)}
    \ControlFlowTok{else}\NormalTok{:}
        \BuiltInTok{print}\NormalTok{(}\StringTok{"x es menor o igual a 5"}\NormalTok{)}
\ControlFlowTok{else}\NormalTok{:}
    \BuiltInTok{print}\NormalTok{(}\StringTok{"x es cero o negativo"}\NormalTok{)}
\end{Highlighting}
\end{Shaded}

\subsection{Ejemplos económicos}\label{ejemplos-econuxf3micos-3}

\begin{Shaded}
\begin{Highlighting}[]
\CommentTok{\# Ejemplo 1: Análisis de elegibilidad para crédito hipotecario}
\NormalTok{ingreso\_mensual }\OperatorTok{=} \DecValTok{8000}
\NormalTok{ahorro\_inicial }\OperatorTok{=} \DecValTok{150000}
\NormalTok{historial\_crediticio }\OperatorTok{=} \StringTok{"bueno"}
\NormalTok{precio\_vivienda }\OperatorTok{=} \DecValTok{500000}

\BuiltInTok{print}\NormalTok{(}\StringTok{"ANÁLISIS DE ELEGIBILIDAD {-} CRÉDITO HIPOTECARIO"}\NormalTok{)}
\BuiltInTok{print}\NormalTok{(}\StringTok{"="} \OperatorTok{*} \DecValTok{60}\NormalTok{)}

\CommentTok{\# Primera condición: ingreso mínimo}
\ControlFlowTok{if}\NormalTok{ ingreso\_mensual }\OperatorTok{\textgreater{}=} \DecValTok{5000}\NormalTok{:}
    \BuiltInTok{print}\NormalTok{(}\SpecialStringTok{f"✓ Ingreso mensual adecuado: S/ }\SpecialCharTok{\{}\NormalTok{ingreso\_mensual}\SpecialCharTok{:,\}}\SpecialStringTok{"}\NormalTok{)}
    
    \CommentTok{\# Segunda condición: ahorro inicial}
\NormalTok{    ahorro\_minimo }\OperatorTok{=}\NormalTok{ precio\_vivienda }\OperatorTok{*} \FloatTok{0.20}
    
    \ControlFlowTok{if}\NormalTok{ ahorro\_inicial }\OperatorTok{\textgreater{}=}\NormalTok{ ahorro\_minimo:}
        \BuiltInTok{print}\NormalTok{(}\SpecialStringTok{f"✓ Ahorro inicial suficiente: S/ }\SpecialCharTok{\{}\NormalTok{ahorro\_inicial}\SpecialCharTok{:,\}}\SpecialStringTok{"}\NormalTok{)}
        \BuiltInTok{print}\NormalTok{(}\SpecialStringTok{f"  (Requerido: S/ }\SpecialCharTok{\{}\NormalTok{ahorro\_minimo}\SpecialCharTok{:,\}}\SpecialStringTok{)"}\NormalTok{)}
        
        \CommentTok{\# Tercera condición: historial crediticio}
        \ControlFlowTok{if}\NormalTok{ historial\_crediticio }\OperatorTok{==} \StringTok{"excelente"}\NormalTok{:}
\NormalTok{            tasa\_interes }\OperatorTok{=} \FloatTok{7.5}
            \BuiltInTok{print}\NormalTok{(}\SpecialStringTok{f"✓ Historial crediticio excelente"}\NormalTok{)}
            \BuiltInTok{print}\NormalTok{(}\SpecialStringTok{f"  Tasa de interés: }\SpecialCharTok{\{}\NormalTok{tasa\_interes}\SpecialCharTok{\}}\SpecialStringTok{\%"}\NormalTok{)}
        \ControlFlowTok{elif}\NormalTok{ historial\_crediticio }\OperatorTok{==} \StringTok{"bueno"}\NormalTok{:}
\NormalTok{            tasa\_interes }\OperatorTok{=} \FloatTok{8.5}
            \BuiltInTok{print}\NormalTok{(}\SpecialStringTok{f"✓ Historial crediticio bueno"}\NormalTok{)}
            \BuiltInTok{print}\NormalTok{(}\SpecialStringTok{f"  Tasa de interés: }\SpecialCharTok{\{}\NormalTok{tasa\_interes}\SpecialCharTok{\}}\SpecialStringTok{\%"}\NormalTok{)}
        \ControlFlowTok{else}\NormalTok{:}
            \BuiltInTok{print}\NormalTok{(}\StringTok{"✗ Historial crediticio insuficiente"}\NormalTok{)}
\NormalTok{            tasa\_interes }\OperatorTok{=} \VariableTok{None}
        
        \CommentTok{\# Decisión final}
        \ControlFlowTok{if}\NormalTok{ tasa\_interes:}
\NormalTok{            monto\_credito }\OperatorTok{=}\NormalTok{ precio\_vivienda }\OperatorTok{{-}}\NormalTok{ ahorro\_inicial}
            \BuiltInTok{print}\NormalTok{(}\SpecialStringTok{f"}\CharTok{\textbackslash{}n}\SpecialCharTok{\{}\StringTok{\textquotesingle{}=\textquotesingle{}}\OperatorTok{*}\DecValTok{60}\SpecialCharTok{\}}\SpecialStringTok{"}\NormalTok{)}
            \BuiltInTok{print}\NormalTok{(}\SpecialStringTok{f"CRÉDITO APROBADO"}\NormalTok{)}
            \BuiltInTok{print}\NormalTok{(}\SpecialStringTok{f"Monto del crédito: S/ }\SpecialCharTok{\{}\NormalTok{monto\_credito}\SpecialCharTok{:,\}}\SpecialStringTok{"}\NormalTok{)}
            \BuiltInTok{print}\NormalTok{(}\SpecialStringTok{f"Tasa de interés: }\SpecialCharTok{\{}\NormalTok{tasa\_interes}\SpecialCharTok{\}}\SpecialStringTok{\%"}\NormalTok{)}
        \ControlFlowTok{else}\NormalTok{:}
            \BuiltInTok{print}\NormalTok{(}\SpecialStringTok{f"}\CharTok{\textbackslash{}n}\SpecialStringTok{Crédito denegado por historial crediticio"}\NormalTok{)}
    \ControlFlowTok{else}\NormalTok{:}
\NormalTok{        faltante }\OperatorTok{=}\NormalTok{ ahorro\_minimo }\OperatorTok{{-}}\NormalTok{ ahorro\_inicial}
        \BuiltInTok{print}\NormalTok{(}\SpecialStringTok{f"✗ Ahorro inicial insuficiente"}\NormalTok{)}
        \BuiltInTok{print}\NormalTok{(}\SpecialStringTok{f"  Ahorro actual: S/ }\SpecialCharTok{\{}\NormalTok{ahorro\_inicial}\SpecialCharTok{:,\}}\SpecialStringTok{"}\NormalTok{)}
        \BuiltInTok{print}\NormalTok{(}\SpecialStringTok{f"  Ahorro requerido: S/ }\SpecialCharTok{\{}\NormalTok{ahorro\_minimo}\SpecialCharTok{:,\}}\SpecialStringTok{"}\NormalTok{)}
        \BuiltInTok{print}\NormalTok{(}\SpecialStringTok{f"  Faltante: S/ }\SpecialCharTok{\{}\NormalTok{faltante}\SpecialCharTok{:,\}}\SpecialStringTok{"}\NormalTok{)}
\ControlFlowTok{else}\NormalTok{:}
    \BuiltInTok{print}\NormalTok{(}\SpecialStringTok{f"✗ Ingreso mensual insuficiente: S/ }\SpecialCharTok{\{}\NormalTok{ingreso\_mensual}\SpecialCharTok{:,\}}\SpecialStringTok{"}\NormalTok{)}
    \BuiltInTok{print}\NormalTok{(}\SpecialStringTok{f"  Ingreso mínimo requerido: S/ 5,000"}\NormalTok{)}

\CommentTok{\# Ejemplo 2: Clasificación de empresa por tamaño y sector}
\NormalTok{ventas\_anuales }\OperatorTok{=} \DecValTok{8500000}  \CommentTok{\# en soles}
\NormalTok{num\_trabajadores }\OperatorTok{=} \DecValTok{85}
\NormalTok{sector }\OperatorTok{=} \StringTok{"manufactura"}

\BuiltInTok{print}\NormalTok{(}\StringTok{"}\CharTok{\textbackslash{}n}\StringTok{"} \OperatorTok{+} \StringTok{"="}\OperatorTok{*}\DecValTok{60}\NormalTok{)}
\BuiltInTok{print}\NormalTok{(}\StringTok{"CLASIFICACIÓN DE EMPRESA"}\NormalTok{)}
\BuiltInTok{print}\NormalTok{(}\StringTok{"="}\OperatorTok{*}\DecValTok{60}\NormalTok{)}

\ControlFlowTok{if}\NormalTok{ sector }\OperatorTok{==} \StringTok{"manufactura"}\NormalTok{:}
    \BuiltInTok{print}\NormalTok{(}\SpecialStringTok{f"Sector: Manufactura"}\NormalTok{)}
    
    \ControlFlowTok{if}\NormalTok{ ventas\_anuales }\OperatorTok{\textless{}=} \DecValTok{1700000}\NormalTok{:}
        \ControlFlowTok{if}\NormalTok{ num\_trabajadores }\OperatorTok{\textless{}=} \DecValTok{10}\NormalTok{:}
\NormalTok{            clasificacion }\OperatorTok{=} \StringTok{"Microempresa"}
        \ControlFlowTok{else}\NormalTok{:}
\NormalTok{            clasificacion }\OperatorTok{=} \StringTok{"Pequeña empresa"}
    \ControlFlowTok{elif}\NormalTok{ ventas\_anuales }\OperatorTok{\textless{}=} \DecValTok{17000000}\NormalTok{:}
        \ControlFlowTok{if}\NormalTok{ num\_trabajadores }\OperatorTok{\textless{}=} \DecValTok{100}\NormalTok{:}
\NormalTok{            clasificacion }\OperatorTok{=} \StringTok{"Pequeña empresa"}
        \ControlFlowTok{else}\NormalTok{:}
\NormalTok{            clasificacion }\OperatorTok{=} \StringTok{"Mediana empresa"}
    \ControlFlowTok{else}\NormalTok{:}
\NormalTok{        clasificacion }\OperatorTok{=} \StringTok{"Gran empresa"}

\ControlFlowTok{elif}\NormalTok{ sector }\OperatorTok{==} \StringTok{"servicios"}\NormalTok{:}
    \BuiltInTok{print}\NormalTok{(}\SpecialStringTok{f"Sector: Servicios"}\NormalTok{)}
    
    \ControlFlowTok{if}\NormalTok{ ventas\_anuales }\OperatorTok{\textless{}=} \DecValTok{1100000}\NormalTok{:}
\NormalTok{        clasificacion }\OperatorTok{=} \StringTok{"Microempresa"}
    \ControlFlowTok{elif}\NormalTok{ ventas\_anuales }\OperatorTok{\textless{}=} \DecValTok{11000000}\NormalTok{:}
\NormalTok{        clasificacion }\OperatorTok{=} \StringTok{"Pequeña empresa"}
    \ControlFlowTok{elif}\NormalTok{ ventas\_anuales }\OperatorTok{\textless{}=} \DecValTok{110000000}\NormalTok{:}
\NormalTok{        clasificacion }\OperatorTok{=} \StringTok{"Mediana empresa"}
    \ControlFlowTok{else}\NormalTok{:}
\NormalTok{        clasificacion }\OperatorTok{=} \StringTok{"Gran empresa"}
\ControlFlowTok{else}\NormalTok{:}
\NormalTok{    clasificacion }\OperatorTok{=} \StringTok{"Sector no clasificado"}

\BuiltInTok{print}\NormalTok{(}\SpecialStringTok{f"Ventas anuales: S/ }\SpecialCharTok{\{}\NormalTok{ventas\_anuales}\SpecialCharTok{:,\}}\SpecialStringTok{"}\NormalTok{)}
\BuiltInTok{print}\NormalTok{(}\SpecialStringTok{f"Número de trabajadores: }\SpecialCharTok{\{}\NormalTok{num\_trabajadores}\SpecialCharTok{\}}\SpecialStringTok{"}\NormalTok{)}
\BuiltInTok{print}\NormalTok{(}\SpecialStringTok{f"Clasificación: }\SpecialCharTok{\{}\NormalTok{clasificacion}\SpecialCharTok{\}}\SpecialStringTok{"}\NormalTok{)}
\end{Highlighting}
\end{Shaded}

\subsection{Alternativa a anidamiento
excesivo}\label{alternativa-a-anidamiento-excesivo}

Los condicionales muy anidados pueden dificultar la lectura. Es mejor
usar \texttt{elif}:

\begin{Shaded}
\begin{Highlighting}[]
\CommentTok{\# Evitar anidamiento excesivo}
\NormalTok{ingreso }\OperatorTok{=} \DecValTok{5000}
\NormalTok{credito }\OperatorTok{=} \StringTok{"bueno"}
\NormalTok{empleo }\OperatorTok{=} \StringTok{"estable"}

\CommentTok{\# Mal: muy anidado}
\ControlFlowTok{if}\NormalTok{ ingreso }\OperatorTok{\textgreater{}} \DecValTok{3000}\NormalTok{:}
    \ControlFlowTok{if}\NormalTok{ credito }\OperatorTok{==} \StringTok{"bueno"}\NormalTok{:}
        \ControlFlowTok{if}\NormalTok{ empleo }\OperatorTok{==} \StringTok{"estable"}\NormalTok{:}
            \BuiltInTok{print}\NormalTok{(}\StringTok{"Crédito aprobado"}\NormalTok{)}
        \ControlFlowTok{else}\NormalTok{:}
            \BuiltInTok{print}\NormalTok{(}\StringTok{"Crédito denegado: empleo inestable"}\NormalTok{)}
    \ControlFlowTok{else}\NormalTok{:}
        \BuiltInTok{print}\NormalTok{(}\StringTok{"Crédito denegado: mal historial"}\NormalTok{)}
\ControlFlowTok{else}\NormalTok{:}
    \BuiltInTok{print}\NormalTok{(}\StringTok{"Crédito denegado: ingreso bajo"}\NormalTok{)}

\CommentTok{\# Mejor: condiciones combinadas}
\ControlFlowTok{if}\NormalTok{ ingreso }\OperatorTok{\textless{}=} \DecValTok{3000}\NormalTok{:}
    \BuiltInTok{print}\NormalTok{(}\StringTok{"Crédito denegado: ingreso bajo"}\NormalTok{)}
\ControlFlowTok{elif}\NormalTok{ credito }\OperatorTok{!=} \StringTok{"bueno"}\NormalTok{:}
    \BuiltInTok{print}\NormalTok{(}\StringTok{"Crédito denegado: mal historial"}\NormalTok{)}
\ControlFlowTok{elif}\NormalTok{ empleo }\OperatorTok{!=} \StringTok{"estable"}\NormalTok{:}
    \BuiltInTok{print}\NormalTok{(}\StringTok{"Crédito denegado: empleo inestable"}\NormalTok{)}
\ControlFlowTok{else}\NormalTok{:}
    \BuiltInTok{print}\NormalTok{(}\StringTok{"Crédito aprobado"}\NormalTok{)}

\CommentTok{\# Aún mejor: usar operadores lógicos}
\ControlFlowTok{if}\NormalTok{ ingreso }\OperatorTok{\textgreater{}} \DecValTok{3000} \KeywordTok{and}\NormalTok{ credito }\OperatorTok{==} \StringTok{"bueno"} \KeywordTok{and}\NormalTok{ empleo }\OperatorTok{==} \StringTok{"estable"}\NormalTok{:}
    \BuiltInTok{print}\NormalTok{(}\StringTok{"Crédito aprobado"}\NormalTok{)}
\ControlFlowTok{else}\NormalTok{:}
    \BuiltInTok{print}\NormalTok{(}\StringTok{"Crédito denegado"}\NormalTok{)}
\end{Highlighting}
\end{Shaded}

\section{Evaluación de pertenencia e
identidad}\label{evaluaciuxf3n-de-pertenencia-e-identidad}

\subsection{Operador in}\label{operador-in}

Verifica si un elemento está en una secuencia:

\begin{Shaded}
\begin{Highlighting}[]
\CommentTok{\# Verificar en listas}
\NormalTok{sectores\_primarios }\OperatorTok{=}\NormalTok{ [}\StringTok{\textquotesingle{}Agricultura\textquotesingle{}}\NormalTok{, }\StringTok{\textquotesingle{}Pesca\textquotesingle{}}\NormalTok{, }\StringTok{\textquotesingle{}Minería\textquotesingle{}}\NormalTok{, }\StringTok{\textquotesingle{}Petróleo\textquotesingle{}}\NormalTok{]}
\NormalTok{sector\_analizar }\OperatorTok{=} \StringTok{\textquotesingle{}Minería\textquotesingle{}}

\ControlFlowTok{if}\NormalTok{ sector\_analizar }\KeywordTok{in}\NormalTok{ sectores\_primarios:}
    \BuiltInTok{print}\NormalTok{(}\SpecialStringTok{f"}\SpecialCharTok{\{}\NormalTok{sector\_analizar}\SpecialCharTok{\}}\SpecialStringTok{ es un sector primario"}\NormalTok{)}
\ControlFlowTok{else}\NormalTok{:}
    \BuiltInTok{print}\NormalTok{(}\SpecialStringTok{f"}\SpecialCharTok{\{}\NormalTok{sector\_analizar}\SpecialCharTok{\}}\SpecialStringTok{ no es un sector primario"}\NormalTok{)}

\CommentTok{\# Verificar en strings}
\NormalTok{descripcion }\OperatorTok{=} \StringTok{"Crecimiento del Producto Bruto Interno"}
\ControlFlowTok{if} \StringTok{"PIB"} \KeywordTok{in}\NormalTok{ descripcion }\KeywordTok{or} \StringTok{"Producto Bruto Interno"} \KeywordTok{in}\NormalTok{ descripcion:}
    \BuiltInTok{print}\NormalTok{(}\StringTok{"El texto menciona el PIB"}\NormalTok{)}

\CommentTok{\# Verificar en diccionarios (verifica claves)}
\NormalTok{indicadores }\OperatorTok{=}\NormalTok{ \{}
    \StringTok{\textquotesingle{}pib\textquotesingle{}}\NormalTok{: }\DecValTok{242632}\NormalTok{,}
    \StringTok{\textquotesingle{}inflacion\textquotesingle{}}\NormalTok{: }\FloatTok{3.5}\NormalTok{,}
    \StringTok{\textquotesingle{}desempleo\textquotesingle{}}\NormalTok{: }\FloatTok{7.2}
\NormalTok{\}}

\ControlFlowTok{if} \StringTok{\textquotesingle{}pib\textquotesingle{}} \KeywordTok{in}\NormalTok{ indicadores:}
    \BuiltInTok{print}\NormalTok{(}\SpecialStringTok{f"PIB disponible: }\SpecialCharTok{\{}\NormalTok{indicadores[}\StringTok{\textquotesingle{}pib\textquotesingle{}}\NormalTok{]}\SpecialCharTok{\}}\SpecialStringTok{"}\NormalTok{)}

\ControlFlowTok{if} \StringTok{\textquotesingle{}reservas\textquotesingle{}} \KeywordTok{not} \KeywordTok{in}\NormalTok{ indicadores:}
    \BuiltInTok{print}\NormalTok{(}\StringTok{"No hay datos de reservas internacionales"}\NormalTok{)}

\CommentTok{\# Verificar en rangos}
\NormalTok{edad }\OperatorTok{=} \DecValTok{25}
\ControlFlowTok{if}\NormalTok{ edad }\KeywordTok{in} \BuiltInTok{range}\NormalTok{(}\DecValTok{18}\NormalTok{, }\DecValTok{65}\NormalTok{):}
    \BuiltInTok{print}\NormalTok{(}\StringTok{"En edad de trabajar"}\NormalTok{)}
\end{Highlighting}
\end{Shaded}

\subsection{Operador is}\label{operador-is}

Verifica identidad de objetos (mismo objeto en memoria):

\begin{Shaded}
\begin{Highlighting}[]
\CommentTok{\# Comparar identidad vs igualdad}
\NormalTok{a }\OperatorTok{=}\NormalTok{ [}\DecValTok{1}\NormalTok{, }\DecValTok{2}\NormalTok{, }\DecValTok{3}\NormalTok{]}
\NormalTok{b }\OperatorTok{=}\NormalTok{ [}\DecValTok{1}\NormalTok{, }\DecValTok{2}\NormalTok{, }\DecValTok{3}\NormalTok{]}
\NormalTok{c }\OperatorTok{=}\NormalTok{ a}

\CommentTok{\# Igualdad de valores}
\BuiltInTok{print}\NormalTok{(a }\OperatorTok{==}\NormalTok{ b)  }\CommentTok{\# True (mismo contenido)}
\BuiltInTok{print}\NormalTok{(a }\OperatorTok{==}\NormalTok{ c)  }\CommentTok{\# True (mismo contenido)}

\CommentTok{\# Identidad de objetos}
\BuiltInTok{print}\NormalTok{(a }\KeywordTok{is}\NormalTok{ b)  }\CommentTok{\# False (objetos diferentes en memoria)}
\BuiltInTok{print}\NormalTok{(a }\KeywordTok{is}\NormalTok{ c)  }\CommentTok{\# True (mismo objeto)}

\CommentTok{\# Uso con None}
\NormalTok{valor }\OperatorTok{=} \VariableTok{None}
\ControlFlowTok{if}\NormalTok{ valor }\KeywordTok{is} \VariableTok{None}\NormalTok{:}
    \BuiltInTok{print}\NormalTok{(}\StringTok{"Valor no definido"}\NormalTok{)}

\CommentTok{\# Con booleanos}
\NormalTok{activo }\OperatorTok{=} \VariableTok{True}
\ControlFlowTok{if}\NormalTok{ activo }\KeywordTok{is} \VariableTok{True}\NormalTok{:}
    \BuiltInTok{print}\NormalTok{(}\StringTok{"Estado activo"}\NormalTok{)}

\CommentTok{\# Mejor práctica: comparar con None usando \textquotesingle{}is\textquotesingle{}}
\NormalTok{datos }\OperatorTok{=} \VariableTok{None}
\ControlFlowTok{if}\NormalTok{ datos }\KeywordTok{is} \VariableTok{None}\NormalTok{:}
    \BuiltInTok{print}\NormalTok{(}\StringTok{"No hay datos disponibles"}\NormalTok{)}
\ControlFlowTok{else}\NormalTok{:}
    \BuiltInTok{print}\NormalTok{(}\SpecialStringTok{f"Datos: }\SpecialCharTok{\{}\NormalTok{datos}\SpecialCharTok{\}}\SpecialStringTok{"}\NormalTok{)}
\end{Highlighting}
\end{Shaded}

\subsection{Función isinstance()}\label{funciuxf3n-isinstance}

Verifica el tipo de un objeto:

\begin{Shaded}
\begin{Highlighting}[]
\CommentTok{\# Verificar tipos}
\NormalTok{pib }\OperatorTok{=} \DecValTok{242632}
\NormalTok{inflacion }\OperatorTok{=} \FloatTok{3.5}
\NormalTok{pais }\OperatorTok{=} \StringTok{"Perú"}
\NormalTok{sectores }\OperatorTok{=}\NormalTok{ [}\StringTok{\textquotesingle{}Minería\textquotesingle{}}\NormalTok{, }\StringTok{\textquotesingle{}Agricultura\textquotesingle{}}\NormalTok{]}
\NormalTok{datos }\OperatorTok{=}\NormalTok{ \{}\StringTok{\textquotesingle{}pib\textquotesingle{}}\NormalTok{: }\DecValTok{242632}\NormalTok{, }\StringTok{\textquotesingle{}inflacion\textquotesingle{}}\NormalTok{: }\FloatTok{3.5}\NormalTok{\}}

\BuiltInTok{print}\NormalTok{(}\BuiltInTok{isinstance}\NormalTok{(pib, }\BuiltInTok{int}\NormalTok{))           }\CommentTok{\# True}
\BuiltInTok{print}\NormalTok{(}\BuiltInTok{isinstance}\NormalTok{(inflacion, }\BuiltInTok{float}\NormalTok{))   }\CommentTok{\# True}
\BuiltInTok{print}\NormalTok{(}\BuiltInTok{isinstance}\NormalTok{(pais, }\BuiltInTok{str}\NormalTok{))          }\CommentTok{\# True}
\BuiltInTok{print}\NormalTok{(}\BuiltInTok{isinstance}\NormalTok{(sectores, }\BuiltInTok{list}\NormalTok{))     }\CommentTok{\# True}
\BuiltInTok{print}\NormalTok{(}\BuiltInTok{isinstance}\NormalTok{(datos, }\BuiltInTok{dict}\NormalTok{))        }\CommentTok{\# True}

\CommentTok{\# Verificar múltiples tipos}
\NormalTok{valor }\OperatorTok{=} \DecValTok{42}
\ControlFlowTok{if} \BuiltInTok{isinstance}\NormalTok{(valor, (}\BuiltInTok{int}\NormalTok{, }\BuiltInTok{float}\NormalTok{)):}
    \BuiltInTok{print}\NormalTok{(}\StringTok{"Es un número"}\NormalTok{)}

\CommentTok{\# Aplicación: validar entrada}
\KeywordTok{def}\NormalTok{ calcular\_interes(capital, tasa, tiempo):}
    \CommentTok{"""Calcula interés simple"""}
    
    \CommentTok{\# Validar tipos}
    \ControlFlowTok{if} \KeywordTok{not} \BuiltInTok{isinstance}\NormalTok{(capital, (}\BuiltInTok{int}\NormalTok{, }\BuiltInTok{float}\NormalTok{)):}
        \ControlFlowTok{return} \StringTok{"Error: capital debe ser numérico"}
    
    \ControlFlowTok{if} \KeywordTok{not} \BuiltInTok{isinstance}\NormalTok{(tasa, (}\BuiltInTok{int}\NormalTok{, }\BuiltInTok{float}\NormalTok{)):}
        \ControlFlowTok{return} \StringTok{"Error: tasa debe ser numérica"}
    
    \ControlFlowTok{if} \KeywordTok{not} \BuiltInTok{isinstance}\NormalTok{(tiempo, (}\BuiltInTok{int}\NormalTok{, }\BuiltInTok{float}\NormalTok{)):}
        \ControlFlowTok{return} \StringTok{"Error: tiempo debe ser numérico"}
    
    \CommentTok{\# Calcular}
\NormalTok{    interes }\OperatorTok{=}\NormalTok{ capital }\OperatorTok{*}\NormalTok{ tasa }\OperatorTok{*}\NormalTok{ tiempo}
    \ControlFlowTok{return}\NormalTok{ interes}

\CommentTok{\# Probar función}
\NormalTok{resultado1 }\OperatorTok{=}\NormalTok{ calcular\_interes(}\DecValTok{10000}\NormalTok{, }\FloatTok{0.08}\NormalTok{, }\DecValTok{3}\NormalTok{)}
\BuiltInTok{print}\NormalTok{(}\SpecialStringTok{f"Interés: S/ }\SpecialCharTok{\{}\NormalTok{resultado1}\SpecialCharTok{:,.2f\}}\SpecialStringTok{"}\NormalTok{)}

\NormalTok{resultado2 }\OperatorTok{=}\NormalTok{ calcular\_interes(}\StringTok{"10000"}\NormalTok{, }\FloatTok{0.08}\NormalTok{, }\DecValTok{3}\NormalTok{)}
\BuiltInTok{print}\NormalTok{(resultado2)  }\CommentTok{\# Error: capital debe ser numérico}
\end{Highlighting}
\end{Shaded}

\section{Guardianes: try y except}\label{guardianes-try-y-except}

El manejo de excepciones permite capturar y manejar errores sin que el
programa se detenga:

\subsection{Estructura básica}\label{estructura-buxe1sica-1}

\begin{Shaded}
\begin{Highlighting}[]
\CommentTok{\# Sin manejo de errores (genera error)}
\CommentTok{\# numero = int("abc")  \# ValueError}

\CommentTok{\# Con manejo de errores}
\ControlFlowTok{try}\NormalTok{:}
\NormalTok{    numero }\OperatorTok{=} \BuiltInTok{int}\NormalTok{(}\StringTok{"abc"}\NormalTok{)}
    \BuiltInTok{print}\NormalTok{(numero)}
\ControlFlowTok{except}\NormalTok{:}
    \BuiltInTok{print}\NormalTok{(}\StringTok{"Error: no se pudo convertir a número"}\NormalTok{)}

\BuiltInTok{print}\NormalTok{(}\StringTok{"El programa continúa"}\NormalTok{)}
\end{Highlighting}
\end{Shaded}

\subsection{Capturar excepciones
específicas}\label{capturar-excepciones-especuxedficas}

\begin{Shaded}
\begin{Highlighting}[]
\CommentTok{\# Capturar tipos específicos de errores}
\ControlFlowTok{try}\NormalTok{:}
\NormalTok{    resultado }\OperatorTok{=} \DecValTok{10} \OperatorTok{/} \DecValTok{0}
\ControlFlowTok{except} \PreprocessorTok{ZeroDivisionError}\NormalTok{:}
    \BuiltInTok{print}\NormalTok{(}\StringTok{"Error: división por cero"}\NormalTok{)}

\ControlFlowTok{try}\NormalTok{:}
\NormalTok{    lista }\OperatorTok{=}\NormalTok{ [}\DecValTok{1}\NormalTok{, }\DecValTok{2}\NormalTok{, }\DecValTok{3}\NormalTok{]}
    \BuiltInTok{print}\NormalTok{(lista[}\DecValTok{5}\NormalTok{])}
\ControlFlowTok{except} \PreprocessorTok{IndexError}\NormalTok{:}
    \BuiltInTok{print}\NormalTok{(}\StringTok{"Error: índice fuera de rango"}\NormalTok{)}

\ControlFlowTok{try}\NormalTok{:}
\NormalTok{    diccionario }\OperatorTok{=}\NormalTok{ \{}\StringTok{\textquotesingle{}a\textquotesingle{}}\NormalTok{: }\DecValTok{1}\NormalTok{, }\StringTok{\textquotesingle{}b\textquotesingle{}}\NormalTok{: }\DecValTok{2}\NormalTok{\}}
    \BuiltInTok{print}\NormalTok{(diccionario[}\StringTok{\textquotesingle{}c\textquotesingle{}}\NormalTok{])}
\ControlFlowTok{except} \PreprocessorTok{KeyError}\NormalTok{:}
    \BuiltInTok{print}\NormalTok{(}\StringTok{"Error: clave no existe"}\NormalTok{)}

\ControlFlowTok{try}\NormalTok{:}
\NormalTok{    numero }\OperatorTok{=} \BuiltInTok{int}\NormalTok{(}\StringTok{"texto"}\NormalTok{)}
\ControlFlowTok{except} \PreprocessorTok{ValueError}\NormalTok{:}
    \BuiltInTok{print}\NormalTok{(}\StringTok{"Error: valor no numérico"}\NormalTok{)}
\end{Highlighting}
\end{Shaded}

\subsection{Múltiples except}\label{muxfaltiples-except}

\begin{Shaded}
\begin{Highlighting}[]
\CommentTok{\# Manejar diferentes errores}
\KeywordTok{def}\NormalTok{ dividir(a, b):}
    \ControlFlowTok{try}\NormalTok{:}
\NormalTok{        resultado }\OperatorTok{=}\NormalTok{ a }\OperatorTok{/}\NormalTok{ b}
        \ControlFlowTok{return}\NormalTok{ resultado}
    \ControlFlowTok{except} \PreprocessorTok{ZeroDivisionError}\NormalTok{:}
        \ControlFlowTok{return} \StringTok{"Error: no se puede dividir por cero"}
    \ControlFlowTok{except} \PreprocessorTok{TypeError}\NormalTok{:}
        \ControlFlowTok{return} \StringTok{"Error: los valores deben ser numéricos"}

\BuiltInTok{print}\NormalTok{(dividir(}\DecValTok{10}\NormalTok{, }\DecValTok{2}\NormalTok{))      }\CommentTok{\# 5.0}
\BuiltInTok{print}\NormalTok{(dividir(}\DecValTok{10}\NormalTok{, }\DecValTok{0}\NormalTok{))      }\CommentTok{\# Error: no se puede dividir por cero}
\BuiltInTok{print}\NormalTok{(dividir(}\DecValTok{10}\NormalTok{, }\StringTok{"2"}\NormalTok{))    }\CommentTok{\# Error: los valores deben ser numéricos}
\end{Highlighting}
\end{Shaded}

\subsection{Cláusula else y finally}\label{cluxe1usula-else-y-finally}

\begin{Shaded}
\begin{Highlighting}[]
\CommentTok{\# else: se ejecuta si no hay error}
\CommentTok{\# finally: siempre se ejecuta}

\ControlFlowTok{try}\NormalTok{:}
\NormalTok{    numero }\OperatorTok{=} \BuiltInTok{int}\NormalTok{(}\BuiltInTok{input}\NormalTok{(}\StringTok{"Ingrese un número: "}\NormalTok{))}
\ControlFlowTok{except} \PreprocessorTok{ValueError}\NormalTok{:}
    \BuiltInTok{print}\NormalTok{(}\StringTok{"Error: debe ingresar un número"}\NormalTok{)}
\ControlFlowTok{else}\NormalTok{:}
    \BuiltInTok{print}\NormalTok{(}\SpecialStringTok{f"Número ingresado: }\SpecialCharTok{\{}\NormalTok{numero}\SpecialCharTok{\}}\SpecialStringTok{"}\NormalTok{)}
\NormalTok{    cuadrado }\OperatorTok{=}\NormalTok{ numero }\OperatorTok{**} \DecValTok{2}
    \BuiltInTok{print}\NormalTok{(}\SpecialStringTok{f"Cuadrado: }\SpecialCharTok{\{}\NormalTok{cuadrado}\SpecialCharTok{\}}\SpecialStringTok{"}\NormalTok{)}
\ControlFlowTok{finally}\NormalTok{:}
    \BuiltInTok{print}\NormalTok{(}\StringTok{"Fin del proceso"}\NormalTok{)}
\end{Highlighting}
\end{Shaded}

\subsection{Aplicaciones económicas}\label{aplicaciones-econuxf3micas}

\begin{Shaded}
\begin{Highlighting}[]
\CommentTok{\# Ejemplo 1: Validar entrada de datos}
\KeywordTok{def}\NormalTok{ calcular\_tasa\_crecimiento(valor\_inicial, valor\_final):}
    \CommentTok{"""Calcula la tasa de crecimiento"""}
    
    \ControlFlowTok{try}\NormalTok{:}
        \CommentTok{\# Intentar convertir a float}
\NormalTok{        v\_inicial }\OperatorTok{=} \BuiltInTok{float}\NormalTok{(valor\_inicial)}
\NormalTok{        v\_final }\OperatorTok{=} \BuiltInTok{float}\NormalTok{(valor\_final)}
        
        \CommentTok{\# Validar valores positivos}
        \ControlFlowTok{if}\NormalTok{ v\_inicial }\OperatorTok{\textless{}=} \DecValTok{0}\NormalTok{:}
            \ControlFlowTok{return} \StringTok{"Error: el valor inicial debe ser positivo"}
        
        \CommentTok{\# Calcular tasa}
\NormalTok{        tasa }\OperatorTok{=}\NormalTok{ ((v\_final }\OperatorTok{{-}}\NormalTok{ v\_inicial) }\OperatorTok{/}\NormalTok{ v\_inicial) }\OperatorTok{*} \DecValTok{100}
        \ControlFlowTok{return}\NormalTok{ tasa}
        
    \ControlFlowTok{except} \PreprocessorTok{ValueError}\NormalTok{:}
        \ControlFlowTok{return} \StringTok{"Error: los valores deben ser numéricos"}
    \ControlFlowTok{except} \PreprocessorTok{ZeroDivisionError}\NormalTok{:}
        \ControlFlowTok{return} \StringTok{"Error: el valor inicial no puede ser cero"}

\CommentTok{\# Probar función}
\BuiltInTok{print}\NormalTok{(calcular\_tasa\_crecimiento(}\DecValTok{100}\NormalTok{, }\DecValTok{110}\NormalTok{))      }\CommentTok{\# 10.0}
\BuiltInTok{print}\NormalTok{(calcular\_tasa\_crecimiento(}\StringTok{"100"}\NormalTok{, }\StringTok{"110"}\NormalTok{))  }\CommentTok{\# 10.0}
\BuiltInTok{print}\NormalTok{(calcular\_tasa\_crecimiento(}\StringTok{"abc"}\NormalTok{, }\DecValTok{110}\NormalTok{))    }\CommentTok{\# Error}
\BuiltInTok{print}\NormalTok{(calcular\_tasa\_crecimiento(}\DecValTok{0}\NormalTok{, }\DecValTok{110}\NormalTok{))        }\CommentTok{\# Error}

\CommentTok{\# Ejemplo 2: Leer datos de diccionario de forma segura}
\NormalTok{datos\_pais }\OperatorTok{=}\NormalTok{ \{}
    \StringTok{\textquotesingle{}nombre\textquotesingle{}}\NormalTok{: }\StringTok{\textquotesingle{}Perú\textquotesingle{}}\NormalTok{,}
    \StringTok{\textquotesingle{}pib\textquotesingle{}}\NormalTok{: }\DecValTok{242632}\NormalTok{,}
    \StringTok{\textquotesingle{}poblacion\textquotesingle{}}\NormalTok{: }\DecValTok{33715471}
\NormalTok{\}}

\KeywordTok{def}\NormalTok{ obtener\_indicador(pais, indicador):}
    \CommentTok{"""Obtiene un indicador del diccionario de forma segura"""}
    
    \ControlFlowTok{try}\NormalTok{:}
\NormalTok{        valor }\OperatorTok{=}\NormalTok{ pais[indicador]}
        \ControlFlowTok{return}\NormalTok{ valor}
    \ControlFlowTok{except} \PreprocessorTok{KeyError}\NormalTok{:}
        \ControlFlowTok{return} \SpecialStringTok{f"Indicador \textquotesingle{}}\SpecialCharTok{\{}\NormalTok{indicador}\SpecialCharTok{\}}\SpecialStringTok{\textquotesingle{} no disponible"}
    \ControlFlowTok{except} \PreprocessorTok{TypeError}\NormalTok{:}
        \ControlFlowTok{return} \StringTok{"Error: formato de datos incorrecto"}

\BuiltInTok{print}\NormalTok{(obtener\_indicador(datos\_pais, }\StringTok{\textquotesingle{}pib\textquotesingle{}}\NormalTok{))         }\CommentTok{\# 242632}
\BuiltInTok{print}\NormalTok{(obtener\_indicador(datos\_pais, }\StringTok{\textquotesingle{}inflacion\textquotesingle{}}\NormalTok{))   }\CommentTok{\# No disponible}

\CommentTok{\# Ejemplo 3: Cálculo con validación}
\KeywordTok{def}\NormalTok{ calcular\_pib\_per\_capita(pib, poblacion):}
    \CommentTok{"""Calcula PIB per cápita con manejo de errores"""}
    
    \ControlFlowTok{try}\NormalTok{:}
\NormalTok{        pib\_pc }\OperatorTok{=}\NormalTok{ (pib }\OperatorTok{/}\NormalTok{ poblacion) }\OperatorTok{*} \DecValTok{1000000}
        \ControlFlowTok{return}\NormalTok{ pib\_pc}
    \ControlFlowTok{except} \PreprocessorTok{ZeroDivisionError}\NormalTok{:}
        \ControlFlowTok{return} \StringTok{"Error: población no puede ser cero"}
    \ControlFlowTok{except} \PreprocessorTok{TypeError}\NormalTok{:}
        \ControlFlowTok{return} \StringTok{"Error: valores deben ser numéricos"}
    \ControlFlowTok{except} \PreprocessorTok{Exception} \ImportTok{as}\NormalTok{ e:}
        \ControlFlowTok{return} \SpecialStringTok{f"Error inesperado: }\SpecialCharTok{\{}\NormalTok{e}\SpecialCharTok{\}}\SpecialStringTok{"}

\BuiltInTok{print}\NormalTok{(calcular\_pib\_per\_capita(}\DecValTok{242632}\NormalTok{, }\DecValTok{33715471}\NormalTok{))    }\CommentTok{\# 7196.21}
\BuiltInTok{print}\NormalTok{(calcular\_pib\_per\_capita(}\DecValTok{242632}\NormalTok{, }\DecValTok{0}\NormalTok{))           }\CommentTok{\# Error: población...}
\BuiltInTok{print}\NormalTok{(calcular\_pib\_per\_capita(}\StringTok{"242632"}\NormalTok{, }\DecValTok{33715471}\NormalTok{))  }\CommentTok{\# Error: valores...}
\end{Highlighting}
\end{Shaded}

\subsection{Obtener información del
error}\label{obtener-informaciuxf3n-del-error}

\begin{Shaded}
\begin{Highlighting}[]
\CommentTok{\# Capturar el mensaje de error}
\ControlFlowTok{try}\NormalTok{:}
\NormalTok{    resultado }\OperatorTok{=} \DecValTok{10} \OperatorTok{/} \DecValTok{0}
\ControlFlowTok{except} \PreprocessorTok{ZeroDivisionError} \ImportTok{as}\NormalTok{ error:}
    \BuiltInTok{print}\NormalTok{(}\SpecialStringTok{f"Error capturado: }\SpecialCharTok{\{}\NormalTok{error}\SpecialCharTok{\}}\SpecialStringTok{"}\NormalTok{)}
    \BuiltInTok{print}\NormalTok{(}\SpecialStringTok{f"Tipo de error: }\SpecialCharTok{\{}\BuiltInTok{type}\NormalTok{(error)}\SpecialCharTok{.}\VariableTok{\_\_name\_\_}\SpecialCharTok{\}}\SpecialStringTok{"}\NormalTok{)}

\CommentTok{\# Aplicación: log de errores}
\KeywordTok{def}\NormalTok{ procesar\_datos(datos):}
\NormalTok{    errores }\OperatorTok{=}\NormalTok{ []}
    
    \ControlFlowTok{for}\NormalTok{ i, valor }\KeywordTok{in} \BuiltInTok{enumerate}\NormalTok{(datos):}
        \ControlFlowTok{try}\NormalTok{:}
\NormalTok{            resultado }\OperatorTok{=} \DecValTok{100} \OperatorTok{/}\NormalTok{ valor}
            \BuiltInTok{print}\NormalTok{(}\SpecialStringTok{f"Índice }\SpecialCharTok{\{}\NormalTok{i}\SpecialCharTok{\}}\SpecialStringTok{: }\SpecialCharTok{\{}\NormalTok{resultado}\SpecialCharTok{:.2f\}}\SpecialStringTok{"}\NormalTok{)}
        \ControlFlowTok{except} \PreprocessorTok{Exception} \ImportTok{as}\NormalTok{ e:}
\NormalTok{            mensaje\_error }\OperatorTok{=} \SpecialStringTok{f"Índice }\SpecialCharTok{\{}\NormalTok{i}\SpecialCharTok{\}}\SpecialStringTok{: }\SpecialCharTok{\{}\BuiltInTok{type}\NormalTok{(e)}\SpecialCharTok{.}\VariableTok{\_\_name\_\_}\SpecialCharTok{\}}\SpecialStringTok{ {-} }\SpecialCharTok{\{}\NormalTok{e}\SpecialCharTok{\}}\SpecialStringTok{"}
\NormalTok{            errores.append(mensaje\_error)}
    
    \ControlFlowTok{if}\NormalTok{ errores:}
        \BuiltInTok{print}\NormalTok{(}\StringTok{"}\CharTok{\textbackslash{}n}\StringTok{Errores encontrados:"}\NormalTok{)}
        \ControlFlowTok{for}\NormalTok{ error }\KeywordTok{in}\NormalTok{ errores:}
            \BuiltInTok{print}\NormalTok{(}\SpecialStringTok{f"  {-} }\SpecialCharTok{\{}\NormalTok{error}\SpecialCharTok{\}}\SpecialStringTok{"}\NormalTok{)}

\NormalTok{datos\_prueba }\OperatorTok{=}\NormalTok{ [}\DecValTok{10}\NormalTok{, }\DecValTok{5}\NormalTok{, }\DecValTok{0}\NormalTok{, }\DecValTok{2}\NormalTok{, }\StringTok{"a"}\NormalTok{, }\DecValTok{4}\NormalTok{]}
\NormalTok{procesar\_datos(datos\_prueba)}
\end{Highlighting}
\end{Shaded}

\section{Evaluación de
cortocircuito}\label{evaluaciuxf3n-de-cortocircuito}

Python evalúa expresiones booleanas de izquierda a derecha y se detiene
cuando el resultado es determinado:

\subsection{Cortocircuito con AND}\label{cortocircuito-con-and}

\begin{Shaded}
\begin{Highlighting}[]
\CommentTok{\# and: se detiene en el primer False}
\KeywordTok{def}\NormalTok{ funcion\_a():}
    \BuiltInTok{print}\NormalTok{(}\StringTok{"Ejecutando función A"}\NormalTok{)}
    \ControlFlowTok{return} \VariableTok{False}

\KeywordTok{def}\NormalTok{ funcion\_b():}
    \BuiltInTok{print}\NormalTok{(}\StringTok{"Ejecutando función B"}\NormalTok{)}
    \ControlFlowTok{return} \VariableTok{True}

\CommentTok{\# funcion\_b nunca se ejecuta}
\NormalTok{resultado }\OperatorTok{=}\NormalTok{ funcion\_a() }\KeywordTok{and}\NormalTok{ funcion\_b()}
\BuiltInTok{print}\NormalTok{(}\SpecialStringTok{f"Resultado: }\SpecialCharTok{\{}\NormalTok{resultado}\SpecialCharTok{\}}\SpecialStringTok{"}\NormalTok{)}
\CommentTok{\# Salida:}
\CommentTok{\# Ejecutando función A}
\CommentTok{\# Resultado: False}

\CommentTok{\# Aplicación: evitar errores}
\NormalTok{datos }\OperatorTok{=}\NormalTok{ []}
\CommentTok{\# Sin cortocircuito, esto daría error:}
\CommentTok{\# if datos[0] \textgreater{} 0:  \# IndexError}

\CommentTok{\# Con cortocircuito, es seguro:}
\ControlFlowTok{if} \BuiltInTok{len}\NormalTok{(datos) }\OperatorTok{\textgreater{}} \DecValTok{0} \KeywordTok{and}\NormalTok{ datos[}\DecValTok{0}\NormalTok{] }\OperatorTok{\textgreater{}} \DecValTok{0}\NormalTok{:}
    \BuiltInTok{print}\NormalTok{(}\StringTok{"Primer elemento es positivo"}\NormalTok{)}
\ControlFlowTok{else}\NormalTok{:}
    \BuiltInTok{print}\NormalTok{(}\StringTok{"Lista vacía o primer elemento no positivo"}\NormalTok{)}
\end{Highlighting}
\end{Shaded}

\subsection{Cortocircuito con OR}\label{cortocircuito-con-or}

\begin{Shaded}
\begin{Highlighting}[]
\CommentTok{\# or: se detiene en el primer True}
\KeywordTok{def}\NormalTok{ funcion\_c():}
    \BuiltInTok{print}\NormalTok{(}\StringTok{"Ejecutando función C"}\NormalTok{)}
    \ControlFlowTok{return} \VariableTok{True}

\KeywordTok{def}\NormalTok{ funcion\_d():}
    \BuiltInTok{print}\NormalTok{(}\StringTok{"Ejecutando función D"}\NormalTok{)}
    \ControlFlowTok{return} \VariableTok{False}

\CommentTok{\# funcion\_d nunca se ejecuta}
\NormalTok{resultado }\OperatorTok{=}\NormalTok{ funcion\_c() }\KeywordTok{or}\NormalTok{ funcion\_d()}
\BuiltInTok{print}\NormalTok{(}\SpecialStringTok{f"Resultado: }\SpecialCharTok{\{}\NormalTok{resultado}\SpecialCharTok{\}}\SpecialStringTok{"}\NormalTok{)}
\CommentTok{\# Salida:}
\CommentTok{\# Ejecutando función C}
\CommentTok{\# Resultado: True}

\CommentTok{\# Aplicación: valores por defecto}
\NormalTok{nombre }\OperatorTok{=} \StringTok{""}
\NormalTok{nombre\_mostrar }\OperatorTok{=}\NormalTok{ nombre }\KeywordTok{or} \StringTok{"Sin nombre"}
\BuiltInTok{print}\NormalTok{(nombre\_mostrar)  }\CommentTok{\# "Sin nombre"}

\NormalTok{pais }\OperatorTok{=} \StringTok{"Perú"}
\NormalTok{pais\_mostrar }\OperatorTok{=}\NormalTok{ pais }\KeywordTok{or} \StringTok{"No especificado"}
\BuiltInTok{print}\NormalTok{(pais\_mostrar)  }\CommentTok{\# "Perú"}
\end{Highlighting}
\end{Shaded}

\subsection{Uso práctico en
economía}\label{uso-pruxe1ctico-en-economuxeda}

\begin{Shaded}
\begin{Highlighting}[]
\CommentTok{\# Ejemplo 1: Validación eficiente}
\KeywordTok{def}\NormalTok{ analizar\_empresa(datos):}
    \CommentTok{"""Analiza datos de empresa con validación eficiente"""}
    
    \CommentTok{\# Verificaciones en orden de eficiencia}
    \CommentTok{\# Se detiene en la primera que falla}
    \ControlFlowTok{if} \KeywordTok{not}\NormalTok{ datos:}
        \ControlFlowTok{return} \StringTok{"Error: no hay datos"}
    
    \ControlFlowTok{if} \StringTok{\textquotesingle{}ingresos\textquotesingle{}} \KeywordTok{not} \KeywordTok{in}\NormalTok{ datos:}
        \ControlFlowTok{return} \StringTok{"Error: faltan ingresos"}
    
    \ControlFlowTok{if} \StringTok{\textquotesingle{}gastos\textquotesingle{}} \KeywordTok{not} \KeywordTok{in}\NormalTok{ datos:}
        \ControlFlowTok{return} \StringTok{"Error: faltan gastos"}
    
    \CommentTok{\# Si llegamos aquí, todos los datos están}
\NormalTok{    utilidad }\OperatorTok{=}\NormalTok{ datos[}\StringTok{\textquotesingle{}ingresos\textquotesingle{}}\NormalTok{] }\OperatorTok{{-}}\NormalTok{ datos[}\StringTok{\textquotesingle{}gastos\textquotesingle{}}\NormalTok{]}
\NormalTok{    margen }\OperatorTok{=}\NormalTok{ (utilidad }\OperatorTok{/}\NormalTok{ datos[}\StringTok{\textquotesingle{}ingresos\textquotesingle{}}\NormalTok{]) }\OperatorTok{*} \DecValTok{100}
    
    \ControlFlowTok{return} \SpecialStringTok{f"Utilidad: S/ }\SpecialCharTok{\{}\NormalTok{utilidad}\SpecialCharTok{:,\}}\SpecialStringTok{, Margen: }\SpecialCharTok{\{}\NormalTok{margen}\SpecialCharTok{:.2f\}}\SpecialStringTok{\%"}

\CommentTok{\# Pruebas}
\BuiltInTok{print}\NormalTok{(analizar\_empresa(\{\}))}
\BuiltInTok{print}\NormalTok{(analizar\_empresa(\{}\StringTok{\textquotesingle{}ingresos\textquotesingle{}}\NormalTok{: }\DecValTok{100000}\NormalTok{\}))}
\BuiltInTok{print}\NormalTok{(analizar\_empresa(\{}\StringTok{\textquotesingle{}ingresos\textquotesingle{}}\NormalTok{: }\DecValTok{100000}\NormalTok{, }\StringTok{\textquotesingle{}gastos\textquotesingle{}}\NormalTok{: }\DecValTok{80000}\NormalTok{\}))}

\CommentTok{\# Ejemplo 2: Acceso seguro a datos anidados}
\NormalTok{pais\_data }\OperatorTok{=}\NormalTok{ \{}
    \StringTok{\textquotesingle{}nombre\textquotesingle{}}\NormalTok{: }\StringTok{\textquotesingle{}Perú\textquotesingle{}}\NormalTok{,}
    \StringTok{\textquotesingle{}economia\textquotesingle{}}\NormalTok{: \{}
        \StringTok{\textquotesingle{}pib\textquotesingle{}}\NormalTok{: }\DecValTok{242632}\NormalTok{,}
        \StringTok{\textquotesingle{}sectores\textquotesingle{}}\NormalTok{: \{}
            \StringTok{\textquotesingle{}mineria\textquotesingle{}}\NormalTok{: }\DecValTok{30000}
\NormalTok{        \}}
\NormalTok{    \}}
\NormalTok{\}}

\CommentTok{\# Acceso con cortocircuito}
\NormalTok{tiene\_mineria }\OperatorTok{=}\NormalTok{ pais\_data }\KeywordTok{and} \OperatorTok{\textbackslash{}}
                \StringTok{\textquotesingle{}economia\textquotesingle{}} \KeywordTok{in}\NormalTok{ pais\_data }\KeywordTok{and} \OperatorTok{\textbackslash{}}
                \StringTok{\textquotesingle{}sectores\textquotesingle{}} \KeywordTok{in}\NormalTok{ pais\_data[}\StringTok{\textquotesingle{}economia\textquotesingle{}}\NormalTok{] }\KeywordTok{and} \OperatorTok{\textbackslash{}}
                \StringTok{\textquotesingle{}mineria\textquotesingle{}} \KeywordTok{in}\NormalTok{ pais\_data[}\StringTok{\textquotesingle{}economia\textquotesingle{}}\NormalTok{][}\StringTok{\textquotesingle{}sectores\textquotesingle{}}\NormalTok{]}

\ControlFlowTok{if}\NormalTok{ tiene\_mineria:}
\NormalTok{    valor\_mineria }\OperatorTok{=}\NormalTok{ pais\_data[}\StringTok{\textquotesingle{}economia\textquotesingle{}}\NormalTok{][}\StringTok{\textquotesingle{}sectores\textquotesingle{}}\NormalTok{][}\StringTok{\textquotesingle{}mineria\textquotesingle{}}\NormalTok{]}
    \BuiltInTok{print}\NormalTok{(}\SpecialStringTok{f"Minería: }\SpecialCharTok{\{}\NormalTok{valor\_mineria}\SpecialCharTok{\}}\SpecialStringTok{"}\NormalTok{)}
\ControlFlowTok{else}\NormalTok{:}
    \BuiltInTok{print}\NormalTok{(}\StringTok{"Datos de minería no disponibles"}\NormalTok{)}

\CommentTok{\# Ejemplo 3: Búsqueda eficiente}
\KeywordTok{def}\NormalTok{ buscar\_indicador(indicadores, nombres\_posibles):}
    \CommentTok{"""Busca un indicador por varios nombres posibles"""}
    
    \CommentTok{\# Se detiene en el primer nombre encontrado}
    \ControlFlowTok{for}\NormalTok{ nombre }\KeywordTok{in}\NormalTok{ nombres\_posibles:}
        \ControlFlowTok{if}\NormalTok{ nombre }\KeywordTok{in}\NormalTok{ indicadores:}
            \ControlFlowTok{return}\NormalTok{ indicadores[nombre]}
    
    \ControlFlowTok{return} \VariableTok{None}

\NormalTok{indicadores }\OperatorTok{=}\NormalTok{ \{}
    \StringTok{\textquotesingle{}producto\_bruto\_interno\textquotesingle{}}\NormalTok{: }\DecValTok{242632}\NormalTok{,}
    \StringTok{\textquotesingle{}tasa\_inflacion\textquotesingle{}}\NormalTok{: }\FloatTok{3.5}\NormalTok{,}
    \StringTok{\textquotesingle{}tasa\_desempleo\textquotesingle{}}\NormalTok{: }\FloatTok{7.2}
\NormalTok{\}}

\CommentTok{\# Buscar PIB con nombres alternativos}
\NormalTok{pib }\OperatorTok{=}\NormalTok{ buscar\_indicador(indicadores, [}\StringTok{\textquotesingle{}pib\textquotesingle{}}\NormalTok{, }\StringTok{\textquotesingle{}producto\_bruto\_interno\textquotesingle{}}\NormalTok{, }\StringTok{\textquotesingle{}gdp\textquotesingle{}}\NormalTok{])}
\BuiltInTok{print}\NormalTok{(}\SpecialStringTok{f"PIB encontrado: }\SpecialCharTok{\{}\NormalTok{pib}\SpecialCharTok{\}}\SpecialStringTok{"}\NormalTok{)}
\end{Highlighting}
\end{Shaded}

\section{Aplicaciones económicas}\label{aplicaciones-econuxf3micas-1}

\subsection{Aplicación 1: Sistema de clasificación de riesgo
país}\label{aplicaciuxf3n-1-sistema-de-clasificaciuxf3n-de-riesgo-pauxeds}

\begin{Shaded}
\begin{Highlighting}[]
\KeywordTok{def}\NormalTok{ clasificar\_riesgo\_pais(datos):}
    \CommentTok{"""}
\CommentTok{    Clasifica el riesgo país según múltiples indicadores}
\CommentTok{    """}
    
    \BuiltInTok{print}\NormalTok{(}\StringTok{"SISTEMA DE CLASIFICACIÓN DE RIESGO PAÍS"}\NormalTok{)}
    \BuiltInTok{print}\NormalTok{(}\StringTok{"="} \OperatorTok{*} \DecValTok{60}\NormalTok{)}
    
    \ControlFlowTok{try}\NormalTok{:}
        \CommentTok{\# Extraer indicadores}
\NormalTok{        pib\_crecimiento }\OperatorTok{=}\NormalTok{ datos.get(}\StringTok{\textquotesingle{}pib\_crecimiento\textquotesingle{}}\NormalTok{, }\DecValTok{0}\NormalTok{)}
\NormalTok{        inflacion }\OperatorTok{=}\NormalTok{ datos.get(}\StringTok{\textquotesingle{}inflacion\textquotesingle{}}\NormalTok{, }\DecValTok{0}\NormalTok{)}
\NormalTok{        deficit\_fiscal }\OperatorTok{=}\NormalTok{ datos.get(}\StringTok{\textquotesingle{}deficit\_fiscal\textquotesingle{}}\NormalTok{, }\DecValTok{0}\NormalTok{)}
\NormalTok{        deuda\_pib }\OperatorTok{=}\NormalTok{ datos.get(}\StringTok{\textquotesingle{}deuda\_pib\textquotesingle{}}\NormalTok{, }\DecValTok{0}\NormalTok{)}
\NormalTok{        reservas\_meses\_importacion }\OperatorTok{=}\NormalTok{ datos.get(}\StringTok{\textquotesingle{}reservas\textquotesingle{}}\NormalTok{, }\DecValTok{0}\NormalTok{)}
        
        \CommentTok{\# Inicializar puntaje}
\NormalTok{        puntaje }\OperatorTok{=} \DecValTok{0}
\NormalTok{        detalles }\OperatorTok{=}\NormalTok{ []}
        
        \CommentTok{\# Evaluar crecimiento económico}
        \ControlFlowTok{if}\NormalTok{ pib\_crecimiento }\OperatorTok{\textgreater{}=} \FloatTok{4.0}\NormalTok{:}
\NormalTok{            puntaje }\OperatorTok{+=} \DecValTok{3}
\NormalTok{            detalles.append(}\StringTok{"✓ Crecimiento sólido (3 puntos)"}\NormalTok{)}
        \ControlFlowTok{elif}\NormalTok{ pib\_crecimiento }\OperatorTok{\textgreater{}=} \FloatTok{2.0}\NormalTok{:}
\NormalTok{            puntaje }\OperatorTok{+=} \DecValTok{2}
\NormalTok{            detalles.append(}\StringTok{"○ Crecimiento moderado (2 puntos)"}\NormalTok{)}
        \ControlFlowTok{elif}\NormalTok{ pib\_crecimiento }\OperatorTok{\textgreater{}=} \DecValTok{0}\NormalTok{:}
\NormalTok{            puntaje }\OperatorTok{+=} \DecValTok{1}
\NormalTok{            detalles.append(}\StringTok{"○ Crecimiento débil (1 punto)"}\NormalTok{)}
        \ControlFlowTok{else}\NormalTok{:}
\NormalTok{            detalles.append(}\StringTok{"✗ Recesión (0 puntos)"}\NormalTok{)}
        
        \CommentTok{\# Evaluar inflación}
        \ControlFlowTok{if}\NormalTok{ inflacion }\OperatorTok{\textless{}=} \FloatTok{3.0}\NormalTok{:}
\NormalTok{            puntaje }\OperatorTok{+=} \DecValTok{3}
\NormalTok{            detalles.append(}\StringTok{"✓ Inflación controlada (3 puntos)"}\NormalTok{)}
        \ControlFlowTok{elif}\NormalTok{ inflacion }\OperatorTok{\textless{}=} \FloatTok{5.0}\NormalTok{:}
\NormalTok{            puntaje }\OperatorTok{+=} \DecValTok{2}
\NormalTok{            detalles.append(}\StringTok{"○ Inflación moderada (2 puntos)"}\NormalTok{)}
        \ControlFlowTok{elif}\NormalTok{ inflacion }\OperatorTok{\textless{}=} \FloatTok{10.0}\NormalTok{:}
\NormalTok{            puntaje }\OperatorTok{+=} \DecValTok{1}
\NormalTok{            detalles.append(}\StringTok{"○ Inflación alta (1 punto)"}\NormalTok{)}
        \ControlFlowTok{else}\NormalTok{:}
\NormalTok{            detalles.append(}\StringTok{"✗ Inflación muy alta (0 puntos)"}\NormalTok{)}
        
        \CommentTok{\# Evaluar balance fiscal}
        \ControlFlowTok{if}\NormalTok{ deficit\_fiscal }\OperatorTok{\textless{}=} \FloatTok{2.0}\NormalTok{:}
\NormalTok{            puntaje }\OperatorTok{+=} \DecValTok{2}
\NormalTok{            detalles.append(}\StringTok{"✓ Balance fiscal saludable (2 puntos)"}\NormalTok{)}
        \ControlFlowTok{elif}\NormalTok{ deficit\_fiscal }\OperatorTok{\textless{}=} \FloatTok{4.0}\NormalTok{:}
\NormalTok{            puntaje }\OperatorTok{+=} \DecValTok{1}
\NormalTok{            detalles.append(}\StringTok{"○ Déficit moderado (1 punto)"}\NormalTok{)}
        \ControlFlowTok{else}\NormalTok{:}
\NormalTok{            detalles.append(}\StringTok{"✗ Déficit elevado (0 puntos)"}\NormalTok{)}
        
        \CommentTok{\# Evaluar deuda pública}
        \ControlFlowTok{if}\NormalTok{ deuda\_pib }\OperatorTok{\textless{}=} \FloatTok{40.0}\NormalTok{:}
\NormalTok{            puntaje }\OperatorTok{+=} \DecValTok{2}
\NormalTok{            detalles.append(}\StringTok{"✓ Deuda baja (2 puntos)"}\NormalTok{)}
        \ControlFlowTok{elif}\NormalTok{ deuda\_pib }\OperatorTok{\textless{}=} \FloatTok{60.0}\NormalTok{:}
\NormalTok{            puntaje }\OperatorTok{+=} \DecValTok{1}
\NormalTok{            detalles.append(}\StringTok{"○ Deuda moderada (1 punto)"}\NormalTok{)}
        \ControlFlowTok{else}\NormalTok{:}
\NormalTok{            detalles.append(}\StringTok{"✗ Deuda alta (0 puntos)"}\NormalTok{)}
        
        \CommentTok{\# Evaluar reservas internacionales}
        \ControlFlowTok{if}\NormalTok{ reservas\_meses\_importacion }\OperatorTok{\textgreater{}=} \FloatTok{6.0}\NormalTok{:}
\NormalTok{            puntaje }\OperatorTok{+=} \DecValTok{2}
\NormalTok{            detalles.append(}\StringTok{"✓ Reservas sólidas (2 puntos)"}\NormalTok{)}
        \ControlFlowTok{elif}\NormalTok{ reservas\_meses\_importacion }\OperatorTok{\textgreater{}=} \FloatTok{3.0}\NormalTok{:}
\NormalTok{            puntaje }\OperatorTok{+=} \DecValTok{1}
\NormalTok{            detalles.append(}\StringTok{"○ Reservas adecuadas (1 punto)"}\NormalTok{)}
        \ControlFlowTok{else}\NormalTok{:}
\NormalTok{            detalles.append(}\StringTok{"✗ Reservas bajas (0 puntos)"}\NormalTok{)}
        
        \CommentTok{\# Determinar clasificación}
        \ControlFlowTok{if}\NormalTok{ puntaje }\OperatorTok{\textgreater{}=} \DecValTok{11}\NormalTok{:}
\NormalTok{            clasificacion }\OperatorTok{=} \StringTok{"AAA"}
\NormalTok{            descripcion }\OperatorTok{=} \StringTok{"Riesgo muy bajo"}
\NormalTok{            spread }\OperatorTok{=} \DecValTok{100}
        \ControlFlowTok{elif}\NormalTok{ puntaje }\OperatorTok{\textgreater{}=} \DecValTok{9}\NormalTok{:}
\NormalTok{            clasificacion }\OperatorTok{=} \StringTok{"AA"}
\NormalTok{            descripcion }\OperatorTok{=} \StringTok{"Riesgo bajo"}
\NormalTok{            spread }\OperatorTok{=} \DecValTok{150}
        \ControlFlowTok{elif}\NormalTok{ puntaje }\OperatorTok{\textgreater{}=} \DecValTok{7}\NormalTok{:}
\NormalTok{            clasificacion }\OperatorTok{=} \StringTok{"A"}
\NormalTok{            descripcion }\OperatorTok{=} \StringTok{"Riesgo moderado"}
\NormalTok{            spread }\OperatorTok{=} \DecValTok{250}
        \ControlFlowTok{elif}\NormalTok{ puntaje }\OperatorTok{\textgreater{}=} \DecValTok{5}\NormalTok{:}
\NormalTok{            clasificacion }\OperatorTok{=} \StringTok{"BBB"}
\NormalTok{            descripcion }\OperatorTok{=} \StringTok{"Riesgo medio"}
\NormalTok{            spread }\OperatorTok{=} \DecValTok{400}
        \ControlFlowTok{elif}\NormalTok{ puntaje }\OperatorTok{\textgreater{}=} \DecValTok{3}\NormalTok{:}
\NormalTok{            clasificacion }\OperatorTok{=} \StringTok{"BB"}
\NormalTok{            descripcion }\OperatorTok{=} \StringTok{"Riesgo alto"}
\NormalTok{            spread }\OperatorTok{=} \DecValTok{600}
        \ControlFlowTok{else}\NormalTok{:}
\NormalTok{            clasificacion }\OperatorTok{=} \StringTok{"B"}
\NormalTok{            descripcion }\OperatorTok{=} \StringTok{"Riesgo muy alto"}
\NormalTok{            spread }\OperatorTok{=} \DecValTok{1000}
        
        \CommentTok{\# Mostrar resultados}
        \BuiltInTok{print}\NormalTok{(}\SpecialStringTok{f"}\CharTok{\textbackslash{}n}\SpecialStringTok{País: }\SpecialCharTok{\{}\NormalTok{datos}\SpecialCharTok{.}\NormalTok{get(}\StringTok{\textquotesingle{}nombre\textquotesingle{}}\NormalTok{, }\StringTok{\textquotesingle{}No especificado\textquotesingle{}}\NormalTok{)}\SpecialCharTok{\}}\SpecialStringTok{"}\NormalTok{)}
        \BuiltInTok{print}\NormalTok{(}\SpecialStringTok{f"}\CharTok{\textbackslash{}n}\SpecialStringTok{Indicadores:"}\NormalTok{)}
        \BuiltInTok{print}\NormalTok{(}\SpecialStringTok{f"  Crecimiento PIB: }\SpecialCharTok{\{}\NormalTok{pib\_crecimiento}\SpecialCharTok{:+.1f\}}\SpecialStringTok{\%"}\NormalTok{)}
        \BuiltInTok{print}\NormalTok{(}\SpecialStringTok{f"  Inflación: }\SpecialCharTok{\{}\NormalTok{inflacion}\SpecialCharTok{:.1f\}}\SpecialStringTok{\%"}\NormalTok{)}
        \BuiltInTok{print}\NormalTok{(}\SpecialStringTok{f"  Déficit fiscal: }\SpecialCharTok{\{}\NormalTok{deficit\_fiscal}\SpecialCharTok{:.1f\}}\SpecialStringTok{\% del PIB"}\NormalTok{)}
        \BuiltInTok{print}\NormalTok{(}\SpecialStringTok{f"  Deuda pública: }\SpecialCharTok{\{}\NormalTok{deuda\_pib}\SpecialCharTok{:.1f\}}\SpecialStringTok{\% del PIB"}\NormalTok{)}
        \BuiltInTok{print}\NormalTok{(}\SpecialStringTok{f"  Reservas: }\SpecialCharTok{\{}\NormalTok{reservas\_meses\_importacion}\SpecialCharTok{:.1f\}}\SpecialStringTok{ meses de importación"}\NormalTok{)}
        
        \BuiltInTok{print}\NormalTok{(}\SpecialStringTok{f"}\CharTok{\textbackslash{}n}\SpecialStringTok{Evaluación:"}\NormalTok{)}
        \ControlFlowTok{for}\NormalTok{ detalle }\KeywordTok{in}\NormalTok{ detalles:}
            \BuiltInTok{print}\NormalTok{(}\SpecialStringTok{f"  }\SpecialCharTok{\{}\NormalTok{detalle}\SpecialCharTok{\}}\SpecialStringTok{"}\NormalTok{)}
        
        \BuiltInTok{print}\NormalTok{(}\SpecialStringTok{f"}\CharTok{\textbackslash{}n}\SpecialCharTok{\{}\StringTok{\textquotesingle{}=\textquotesingle{}}\OperatorTok{*}\DecValTok{60}\SpecialCharTok{\}}\SpecialStringTok{"}\NormalTok{)}
        \BuiltInTok{print}\NormalTok{(}\SpecialStringTok{f"Puntaje total: }\SpecialCharTok{\{}\NormalTok{puntaje}\SpecialCharTok{\}}\SpecialStringTok{/12"}\NormalTok{)}
        \BuiltInTok{print}\NormalTok{(}\SpecialStringTok{f"Clasificación: }\SpecialCharTok{\{}\NormalTok{clasificacion}\SpecialCharTok{\}}\SpecialStringTok{"}\NormalTok{)}
        \BuiltInTok{print}\NormalTok{(}\SpecialStringTok{f"Descripción: }\SpecialCharTok{\{}\NormalTok{descripcion}\SpecialCharTok{\}}\SpecialStringTok{"}\NormalTok{)}
        \BuiltInTok{print}\NormalTok{(}\SpecialStringTok{f"Spread estimado: }\SpecialCharTok{\{}\NormalTok{spread}\SpecialCharTok{\}}\SpecialStringTok{ puntos básicos"}\NormalTok{)}
        \BuiltInTok{print}\NormalTok{(}\SpecialStringTok{f"}\SpecialCharTok{\{}\StringTok{\textquotesingle{}=\textquotesingle{}}\OperatorTok{*}\DecValTok{60}\SpecialCharTok{\}}\SpecialStringTok{"}\NormalTok{)}
        
        \ControlFlowTok{return}\NormalTok{ clasificacion}
        
    \ControlFlowTok{except} \PreprocessorTok{Exception} \ImportTok{as}\NormalTok{ e:}
        \ControlFlowTok{return} \SpecialStringTok{f"Error en el análisis: }\SpecialCharTok{\{}\NormalTok{e}\SpecialCharTok{\}}\SpecialStringTok{"}

\CommentTok{\# Ejemplo de uso}
\NormalTok{peru\_datos }\OperatorTok{=}\NormalTok{ \{}
    \StringTok{\textquotesingle{}nombre\textquotesingle{}}\NormalTok{: }\StringTok{\textquotesingle{}Perú\textquotesingle{}}\NormalTok{,}
    \StringTok{\textquotesingle{}pib\_crecimiento\textquotesingle{}}\NormalTok{: }\FloatTok{2.7}\NormalTok{,}
    \StringTok{\textquotesingle{}inflacion\textquotesingle{}}\NormalTok{: }\FloatTok{3.5}\NormalTok{,}
    \StringTok{\textquotesingle{}deficit\_fiscal\textquotesingle{}}\NormalTok{: }\FloatTok{2.8}\NormalTok{,}
    \StringTok{\textquotesingle{}deuda\_pib\textquotesingle{}}\NormalTok{: }\FloatTok{35.0}\NormalTok{,}
    \StringTok{\textquotesingle{}reservas\textquotesingle{}}\NormalTok{: }\FloatTok{8.5}
\NormalTok{\}}

\NormalTok{clasificacion }\OperatorTok{=}\NormalTok{ clasificar\_riesgo\_pais(peru\_datos)}

\CommentTok{\# Comparar con otro país}
\NormalTok{venezuela\_datos }\OperatorTok{=}\NormalTok{ \{}
    \StringTok{\textquotesingle{}nombre\textquotesingle{}}\NormalTok{: }\StringTok{\textquotesingle{}Venezuela\textquotesingle{}}\NormalTok{,}
    \StringTok{\textquotesingle{}pib\_crecimiento\textquotesingle{}}\NormalTok{: }\OperatorTok{{-}}\FloatTok{5.0}\NormalTok{,}
    \StringTok{\textquotesingle{}inflacion\textquotesingle{}}\NormalTok{: }\FloatTok{1500.0}\NormalTok{,}
    \StringTok{\textquotesingle{}deficit\_fiscal\textquotesingle{}}\NormalTok{: }\FloatTok{15.0}\NormalTok{,}
    \StringTok{\textquotesingle{}deuda\_pib\textquotesingle{}}\NormalTok{: }\FloatTok{150.0}\NormalTok{,}
    \StringTok{\textquotesingle{}reservas\textquotesingle{}}\NormalTok{: }\FloatTok{0.5}
\NormalTok{\}}

\BuiltInTok{print}\NormalTok{(}\StringTok{"}\CharTok{\textbackslash{}n}\StringTok{"}\NormalTok{)}
\NormalTok{clasificacion\_vzla }\OperatorTok{=}\NormalTok{ clasificar\_riesgo\_pais(venezuela\_datos)}
\end{Highlighting}
\end{Shaded}

\subsection{Aplicación 2: Calculadora de impuestos
progresivos}\label{aplicaciuxf3n-2-calculadora-de-impuestos-progresivos}

\begin{Shaded}
\begin{Highlighting}[]
\KeywordTok{def}\NormalTok{ calcular\_impuesto\_renta(ingreso\_anual):}
    \CommentTok{"""}
\CommentTok{    Calcula el impuesto a la renta con sistema progresivo}
\CommentTok{    (Basado en el sistema peruano simplificado)}
\CommentTok{    """}
    
    \BuiltInTok{print}\NormalTok{(}\StringTok{"CÁLCULO DE IMPUESTO A LA RENTA"}\NormalTok{)}
    \BuiltInTok{print}\NormalTok{(}\StringTok{"="} \OperatorTok{*} \DecValTok{60}\NormalTok{)}
    
    \ControlFlowTok{try}\NormalTok{:}
        \CommentTok{\# Validar entrada}
        \ControlFlowTok{if} \KeywordTok{not} \BuiltInTok{isinstance}\NormalTok{(ingreso\_anual, (}\BuiltInTok{int}\NormalTok{, }\BuiltInTok{float}\NormalTok{)):}
            \ControlFlowTok{return} \StringTok{"Error: el ingreso debe ser numérico"}
        
        \ControlFlowTok{if}\NormalTok{ ingreso\_anual }\OperatorTok{\textless{}} \DecValTok{0}\NormalTok{:}
            \ControlFlowTok{return} \StringTok{"Error: el ingreso no puede ser negativo"}
        
        \CommentTok{\# Tramos de impuesto (valores en UIT)}
\NormalTok{        UIT }\OperatorTok{=} \DecValTok{5150}  \CommentTok{\# UIT para 2026 (valor hipotético)}
        
        \CommentTok{\# Definir tramos}
\NormalTok{        tramo1\_limite }\OperatorTok{=} \DecValTok{5} \OperatorTok{*}\NormalTok{ UIT      }\CommentTok{\# 0{-}5 UIT: 8\%}
\NormalTok{        tramo2\_limite }\OperatorTok{=} \DecValTok{20} \OperatorTok{*}\NormalTok{ UIT     }\CommentTok{\# 5{-}20 UIT: 14\%}
\NormalTok{        tramo3\_limite }\OperatorTok{=} \DecValTok{35} \OperatorTok{*}\NormalTok{ UIT     }\CommentTok{\# 20{-}35 UIT: 17\%}
\NormalTok{        tramo4\_limite }\OperatorTok{=} \DecValTok{45} \OperatorTok{*}\NormalTok{ UIT     }\CommentTok{\# 35{-}45 UIT: 20\%}
        \CommentTok{\# Más de 45 UIT: 30\%}
        
        \CommentTok{\# Calcular impuesto}
\NormalTok{        impuesto }\OperatorTok{=} \DecValTok{0}
\NormalTok{        ingreso\_restante }\OperatorTok{=}\NormalTok{ ingreso\_anual}
\NormalTok{        desglose }\OperatorTok{=}\NormalTok{ []}
        
        \CommentTok{\# Tramo 1: 0{-}5 UIT (8\%)}
        \ControlFlowTok{if}\NormalTok{ ingreso\_restante }\OperatorTok{\textgreater{}} \DecValTok{0}\NormalTok{:}
\NormalTok{            monto\_tramo }\OperatorTok{=} \BuiltInTok{min}\NormalTok{(ingreso\_restante, tramo1\_limite)}
\NormalTok{            impuesto\_tramo }\OperatorTok{=}\NormalTok{ monto\_tramo }\OperatorTok{*} \FloatTok{0.08}
\NormalTok{            impuesto }\OperatorTok{+=}\NormalTok{ impuesto\_tramo}
\NormalTok{            ingreso\_restante }\OperatorTok{{-}=}\NormalTok{ monto\_tramo}
            
            \ControlFlowTok{if}\NormalTok{ monto\_tramo }\OperatorTok{\textgreater{}} \DecValTok{0}\NormalTok{:}
\NormalTok{                desglose.append(\{}
                    \StringTok{\textquotesingle{}tramo\textquotesingle{}}\NormalTok{: }\StringTok{\textquotesingle{}Tramo 1 (0{-}5 UIT)\textquotesingle{}}\NormalTok{,}
                    \StringTok{\textquotesingle{}base\textquotesingle{}}\NormalTok{: monto\_tramo,}
                    \StringTok{\textquotesingle{}tasa\textquotesingle{}}\NormalTok{: }\FloatTok{0.08}\NormalTok{,}
                    \StringTok{\textquotesingle{}impuesto\textquotesingle{}}\NormalTok{: impuesto\_tramo}
\NormalTok{                \})}
        
        \CommentTok{\# Tramo 2: 5{-}20 UIT (14\%)}
        \ControlFlowTok{if}\NormalTok{ ingreso\_restante }\OperatorTok{\textgreater{}} \DecValTok{0}\NormalTok{:}
\NormalTok{            monto\_tramo }\OperatorTok{=} \BuiltInTok{min}\NormalTok{(ingreso\_restante, tramo2\_limite }\OperatorTok{{-}}\NormalTok{ tramo1\_limite)}
\NormalTok{            impuesto\_tramo }\OperatorTok{=}\NormalTok{ monto\_tramo }\OperatorTok{*} \FloatTok{0.14}
\NormalTok{            impuesto }\OperatorTok{+=}\NormalTok{ impuesto\_tramo}
\NormalTok{            ingreso\_restante }\OperatorTok{{-}=}\NormalTok{ monto\_tramo}
            
            \ControlFlowTok{if}\NormalTok{ monto\_tramo }\OperatorTok{\textgreater{}} \DecValTok{0}\NormalTok{:}
\NormalTok{                desglose.append(\{}
                    \StringTok{\textquotesingle{}tramo\textquotesingle{}}\NormalTok{: }\StringTok{\textquotesingle{}Tramo 2 (5{-}20 UIT)\textquotesingle{}}\NormalTok{,}
                    \StringTok{\textquotesingle{}base\textquotesingle{}}\NormalTok{: monto\_tramo,}
                    \StringTok{\textquotesingle{}tasa\textquotesingle{}}\NormalTok{: }\FloatTok{0.14}\NormalTok{,}
                    \StringTok{\textquotesingle{}impuesto\textquotesingle{}}\NormalTok{: impuesto\_tramo}
\NormalTok{                \})}
        
        \CommentTok{\# Tramo 3: 20{-}35 UIT (17\%)}
        \ControlFlowTok{if}\NormalTok{ ingreso\_restante }\OperatorTok{\textgreater{}} \DecValTok{0}\NormalTok{:}
\NormalTok{            monto\_tramo }\OperatorTok{=} \BuiltInTok{min}\NormalTok{(ingreso\_restante, tramo3\_limite }\OperatorTok{{-}}\NormalTok{ tramo2\_limite)}
\NormalTok{            impuesto\_tramo }\OperatorTok{=}\NormalTok{ monto\_tramo }\OperatorTok{*} \FloatTok{0.17}
\NormalTok{            impuesto }\OperatorTok{+=}\NormalTok{ impuesto\_tramo}
\NormalTok{            ingreso\_restante }\OperatorTok{{-}=}\NormalTok{ monto\_tramo}
            
            \ControlFlowTok{if}\NormalTok{ monto\_tramo }\OperatorTok{\textgreater{}} \DecValTok{0}\NormalTok{:}
\NormalTok{                desglose.append(\{}
                    \StringTok{\textquotesingle{}tramo\textquotesingle{}}\NormalTok{: }\StringTok{\textquotesingle{}Tramo 3 (20{-}35 UIT)\textquotesingle{}}\NormalTok{,}
                    \StringTok{\textquotesingle{}base\textquotesingle{}}\NormalTok{: monto\_tramo,}
                    \StringTok{\textquotesingle{}tasa\textquotesingle{}}\NormalTok{: }\FloatTok{0.17}\NormalTok{,}
                    \StringTok{\textquotesingle{}impuesto\textquotesingle{}}\NormalTok{: impuesto\_tramo}
\NormalTok{                \})}
        
        \CommentTok{\# Tramo 4: 35{-}45 UIT (20\%)}
        \ControlFlowTok{if}\NormalTok{ ingreso\_restante }\OperatorTok{\textgreater{}} \DecValTok{0}\NormalTok{:}
\NormalTok{            monto\_tramo }\OperatorTok{=} \BuiltInTok{min}\NormalTok{(ingreso\_restante, tramo4\_limite }\OperatorTok{{-}}\NormalTok{ tramo3\_limite)}
\NormalTok{            impuesto\_tramo }\OperatorTok{=}\NormalTok{ monto\_tramo }\OperatorTok{*} \FloatTok{0.20}
\NormalTok{            impuesto }\OperatorTok{+=}\NormalTok{ impuesto\_tramo}
\NormalTok{            ingreso\_restante }\OperatorTok{{-}=}\NormalTok{ monto\_tramo}
            
            \ControlFlowTok{if}\NormalTok{ monto\_tramo }\OperatorTok{\textgreater{}} \DecValTok{0}\NormalTok{:}
\NormalTok{                desglose.append(\{}
                    \StringTok{\textquotesingle{}tramo\textquotesingle{}}\NormalTok{: }\StringTok{\textquotesingle{}Tramo 4 (35{-}45 UIT)\textquotesingle{}}\NormalTok{,}
                    \StringTok{\textquotesingle{}base\textquotesingle{}}\NormalTok{: monto\_tramo,}
                    \StringTok{\textquotesingle{}tasa\textquotesingle{}}\NormalTok{: }\FloatTok{0.20}\NormalTok{,}
                    \StringTok{\textquotesingle{}impuesto\textquotesingle{}}\NormalTok{: impuesto\_tramo}
\NormalTok{                \})}
        
        \CommentTok{\# Tramo 5: Más de 45 UIT (30\%)}
        \ControlFlowTok{if}\NormalTok{ ingreso\_restante }\OperatorTok{\textgreater{}} \DecValTok{0}\NormalTok{:}
\NormalTok{            monto\_tramo }\OperatorTok{=}\NormalTok{ ingreso\_restante}
\NormalTok{            impuesto\_tramo }\OperatorTok{=}\NormalTok{ monto\_tramo }\OperatorTok{*} \FloatTok{0.30}
\NormalTok{            impuesto }\OperatorTok{+=}\NormalTok{ impuesto\_tramo}
            
\NormalTok{            desglose.append(\{}
                \StringTok{\textquotesingle{}tramo\textquotesingle{}}\NormalTok{: }\StringTok{\textquotesingle{}Tramo 5 (\textgreater{}45 UIT)\textquotesingle{}}\NormalTok{,}
                \StringTok{\textquotesingle{}base\textquotesingle{}}\NormalTok{: monto\_tramo,}
                \StringTok{\textquotesingle{}tasa\textquotesingle{}}\NormalTok{: }\FloatTok{0.30}\NormalTok{,}
                \StringTok{\textquotesingle{}impuesto\textquotesingle{}}\NormalTok{: impuesto\_tramo}
\NormalTok{            \})}
        
        \CommentTok{\# Calcular tasa efectiva}
\NormalTok{        tasa\_efectiva }\OperatorTok{=}\NormalTok{ (impuesto }\OperatorTok{/}\NormalTok{ ingreso\_anual) }\OperatorTok{*} \DecValTok{100} \ControlFlowTok{if}\NormalTok{ ingreso\_anual }\OperatorTok{\textgreater{}} \DecValTok{0} \ControlFlowTok{else} \DecValTok{0}
\NormalTok{        ingreso\_neto }\OperatorTok{=}\NormalTok{ ingreso\_anual }\OperatorTok{{-}}\NormalTok{ impuesto}
        
        \CommentTok{\# Mostrar resultados}
        \BuiltInTok{print}\NormalTok{(}\SpecialStringTok{f"}\CharTok{\textbackslash{}n}\SpecialStringTok{Ingreso bruto anual: S/ }\SpecialCharTok{\{}\NormalTok{ingreso\_anual}\SpecialCharTok{:,.2f\}}\SpecialStringTok{"}\NormalTok{)}
        \BuiltInTok{print}\NormalTok{(}\SpecialStringTok{f"UIT vigente: S/ }\SpecialCharTok{\{}\NormalTok{UIT}\SpecialCharTok{:,.2f\}}\SpecialStringTok{"}\NormalTok{)}
        \BuiltInTok{print}\NormalTok{(}\SpecialStringTok{f"}\CharTok{\textbackslash{}n}\SpecialStringTok{Desglose por tramos:"}\NormalTok{)}
        \BuiltInTok{print}\NormalTok{(}\StringTok{"{-}"} \OperatorTok{*} \DecValTok{60}\NormalTok{)}
        
        \ControlFlowTok{for}\NormalTok{ item }\KeywordTok{in}\NormalTok{ desglose:}
            \BuiltInTok{print}\NormalTok{(}\SpecialStringTok{f"}\CharTok{\textbackslash{}n}\SpecialCharTok{\{}\NormalTok{item[}\StringTok{\textquotesingle{}tramo\textquotesingle{}}\NormalTok{]}\SpecialCharTok{\}}\SpecialStringTok{"}\NormalTok{)}
            \BuiltInTok{print}\NormalTok{(}\SpecialStringTok{f"  Base imponible: S/ }\SpecialCharTok{\{}\NormalTok{item[}\StringTok{\textquotesingle{}base\textquotesingle{}}\NormalTok{]}\SpecialCharTok{:,.2f\}}\SpecialStringTok{"}\NormalTok{)}
            \BuiltInTok{print}\NormalTok{(}\SpecialStringTok{f"  Tasa: }\SpecialCharTok{\{}\NormalTok{item[}\StringTok{\textquotesingle{}tasa\textquotesingle{}}\NormalTok{]}\OperatorTok{*}\DecValTok{100}\SpecialCharTok{:.0f\}}\SpecialStringTok{\%"}\NormalTok{)}
            \BuiltInTok{print}\NormalTok{(}\SpecialStringTok{f"  Impuesto: S/ }\SpecialCharTok{\{}\NormalTok{item[}\StringTok{\textquotesingle{}impuesto\textquotesingle{}}\NormalTok{]}\SpecialCharTok{:,.2f\}}\SpecialStringTok{"}\NormalTok{)}
        
        \BuiltInTok{print}\NormalTok{(}\SpecialStringTok{f"}\CharTok{\textbackslash{}n}\SpecialCharTok{\{}\StringTok{\textquotesingle{}=\textquotesingle{}}\OperatorTok{*}\DecValTok{60}\SpecialCharTok{\}}\SpecialStringTok{"}\NormalTok{)}
        \BuiltInTok{print}\NormalTok{(}\SpecialStringTok{f"Impuesto total: S/ }\SpecialCharTok{\{}\NormalTok{impuesto}\SpecialCharTok{:,.2f\}}\SpecialStringTok{"}\NormalTok{)}
        \BuiltInTok{print}\NormalTok{(}\SpecialStringTok{f"Tasa efectiva: }\SpecialCharTok{\{}\NormalTok{tasa\_efectiva}\SpecialCharTok{:.2f\}}\SpecialStringTok{\%"}\NormalTok{)}
        \BuiltInTok{print}\NormalTok{(}\SpecialStringTok{f"Ingreso neto: S/ }\SpecialCharTok{\{}\NormalTok{ingreso\_neto}\SpecialCharTok{:,.2f\}}\SpecialStringTok{"}\NormalTok{)}
        \BuiltInTok{print}\NormalTok{(}\SpecialStringTok{f"}\SpecialCharTok{\{}\StringTok{\textquotesingle{}=\textquotesingle{}}\OperatorTok{*}\DecValTok{60}\SpecialCharTok{\}}\SpecialStringTok{"}\NormalTok{)}
        
        \ControlFlowTok{return}\NormalTok{ impuesto}
        
    \ControlFlowTok{except} \PreprocessorTok{Exception} \ImportTok{as}\NormalTok{ e:}
        \ControlFlowTok{return} \SpecialStringTok{f"Error en el cálculo: }\SpecialCharTok{\{}\NormalTok{e}\SpecialCharTok{\}}\SpecialStringTok{"}

\CommentTok{\# Ejemplos de uso}
\BuiltInTok{print}\NormalTok{(}\StringTok{"Caso 1: Ingreso medio"}\NormalTok{)}
\NormalTok{calcular\_impuesto\_renta(}\DecValTok{80000}\NormalTok{)}

\BuiltInTok{print}\NormalTok{(}\StringTok{"}\CharTok{\textbackslash{}n}\StringTok{"} \OperatorTok{+} \StringTok{"="}\OperatorTok{*}\DecValTok{80} \OperatorTok{+} \StringTok{"}\CharTok{\textbackslash{}n}\StringTok{"}\NormalTok{)}

\BuiltInTok{print}\NormalTok{(}\StringTok{"Caso 2: Ingreso alto"}\NormalTok{)}
\NormalTok{calcular\_impuesto\_renta(}\DecValTok{300000}\NormalTok{)}
\end{Highlighting}
\end{Shaded}

\section{Ejercicios prácticos}\label{ejercicios-pruxe1cticos}

\subsection{Ejercicio 1: Sistema de evaluación de proyectos de
inversión}\label{ejercicio-1-sistema-de-evaluaciuxf3n-de-proyectos-de-inversiuxf3n}

\begin{Shaded}
\begin{Highlighting}[]
\KeywordTok{def}\NormalTok{ evaluar\_proyecto(datos\_proyecto):}
    \CommentTok{"""}
\CommentTok{    Evalúa la viabilidad de un proyecto de inversión}
\CommentTok{    """}
    
    \BuiltInTok{print}\NormalTok{(}\StringTok{"EVALUACIÓN DE PROYECTO DE INVERSIÓN"}\NormalTok{)}
    \BuiltInTok{print}\NormalTok{(}\StringTok{"="} \OperatorTok{*} \DecValTok{60}\NormalTok{)}
    
    \ControlFlowTok{try}\NormalTok{:}
        \CommentTok{\# Extraer datos}
\NormalTok{        inversion\_inicial }\OperatorTok{=}\NormalTok{ datos\_proyecto[}\StringTok{\textquotesingle{}inversion\_inicial\textquotesingle{}}\NormalTok{]}
\NormalTok{        flujos\_anuales }\OperatorTok{=}\NormalTok{ datos\_proyecto[}\StringTok{\textquotesingle{}flujos\_anuales\textquotesingle{}}\NormalTok{]}
\NormalTok{        tasa\_descuento }\OperatorTok{=}\NormalTok{ datos\_proyecto[}\StringTok{\textquotesingle{}tasa\_descuento\textquotesingle{}}\NormalTok{]}
\NormalTok{        nombre }\OperatorTok{=}\NormalTok{ datos\_proyecto.get(}\StringTok{\textquotesingle{}nombre\textquotesingle{}}\NormalTok{, }\StringTok{\textquotesingle{}Proyecto sin nombre\textquotesingle{}}\NormalTok{)}
        
        \BuiltInTok{print}\NormalTok{(}\SpecialStringTok{f"}\CharTok{\textbackslash{}n}\SpecialStringTok{Proyecto: }\SpecialCharTok{\{}\NormalTok{nombre}\SpecialCharTok{\}}\SpecialStringTok{"}\NormalTok{)}
        \BuiltInTok{print}\NormalTok{(}\SpecialStringTok{f"Inversión inicial: S/ }\SpecialCharTok{\{}\NormalTok{inversion\_inicial}\SpecialCharTok{:,.2f\}}\SpecialStringTok{"}\NormalTok{)}
        \BuiltInTok{print}\NormalTok{(}\SpecialStringTok{f"Tasa de descuento: }\SpecialCharTok{\{}\NormalTok{tasa\_descuento}\OperatorTok{*}\DecValTok{100}\SpecialCharTok{:.1f\}}\SpecialStringTok{\%"}\NormalTok{)}
        
        \CommentTok{\# Calcular VPN}
\NormalTok{        vpn }\OperatorTok{=} \OperatorTok{{-}}\NormalTok{inversion\_inicial}
        \BuiltInTok{print}\NormalTok{(}\SpecialStringTok{f"}\CharTok{\textbackslash{}n}\SpecialStringTok{Flujos de caja:"}\NormalTok{)}
        
        \ControlFlowTok{for}\NormalTok{ i, flujo }\KeywordTok{in} \BuiltInTok{enumerate}\NormalTok{(flujos\_anuales, }\DecValTok{1}\NormalTok{):}
\NormalTok{            valor\_presente }\OperatorTok{=}\NormalTok{ flujo }\OperatorTok{/}\NormalTok{ ((}\DecValTok{1} \OperatorTok{+}\NormalTok{ tasa\_descuento) }\OperatorTok{**}\NormalTok{ i)}
\NormalTok{            vpn }\OperatorTok{+=}\NormalTok{ valor\_presente}
            \BuiltInTok{print}\NormalTok{(}\SpecialStringTok{f"  Año }\SpecialCharTok{\{}\NormalTok{i}\SpecialCharTok{\}}\SpecialStringTok{: S/ }\SpecialCharTok{\{}\NormalTok{flujo}\SpecialCharTok{:,.2f\}}\SpecialStringTok{ (VP: S/ }\SpecialCharTok{\{}\NormalTok{valor\_presente}\SpecialCharTok{:,.2f\}}\SpecialStringTok{)"}\NormalTok{)}
        
        \CommentTok{\# Calcular TIR (aproximación simple)}
        \CommentTok{\# Para cálculo exacto se requeriría método numérico}
        
        \CommentTok{\# Período de recuperación}
\NormalTok{        saldo\_acumulado }\OperatorTok{=} \OperatorTok{{-}}\NormalTok{inversion\_inicial}
\NormalTok{        periodo\_recuperacion }\OperatorTok{=} \VariableTok{None}
        
        \BuiltInTok{print}\NormalTok{(}\SpecialStringTok{f"}\CharTok{\textbackslash{}n}\SpecialStringTok{Análisis de recuperación:"}\NormalTok{)}
        \ControlFlowTok{for}\NormalTok{ i, flujo }\KeywordTok{in} \BuiltInTok{enumerate}\NormalTok{(flujos\_anuales, }\DecValTok{1}\NormalTok{):}
\NormalTok{            saldo\_acumulado }\OperatorTok{+=}\NormalTok{ flujo}
            \BuiltInTok{print}\NormalTok{(}\SpecialStringTok{f"  Año }\SpecialCharTok{\{}\NormalTok{i}\SpecialCharTok{\}}\SpecialStringTok{: Saldo acumulado S/ }\SpecialCharTok{\{}\NormalTok{saldo\_acumulado}\SpecialCharTok{:,.2f\}}\SpecialStringTok{"}\NormalTok{)}
            
            \ControlFlowTok{if}\NormalTok{ saldo\_acumulado }\OperatorTok{\textgreater{}=} \DecValTok{0} \KeywordTok{and}\NormalTok{ periodo\_recuperacion }\KeywordTok{is} \VariableTok{None}\NormalTok{:}
\NormalTok{                periodo\_recuperacion }\OperatorTok{=}\NormalTok{ i}
        
        \CommentTok{\# Evaluación}
        \BuiltInTok{print}\NormalTok{(}\SpecialStringTok{f"}\CharTok{\textbackslash{}n}\SpecialCharTok{\{}\StringTok{\textquotesingle{}=\textquotesingle{}}\OperatorTok{*}\DecValTok{60}\SpecialCharTok{\}}\SpecialStringTok{"}\NormalTok{)}
        \BuiltInTok{print}\NormalTok{(}\SpecialStringTok{f"RESULTADOS:"}\NormalTok{)}
        \BuiltInTok{print}\NormalTok{(}\SpecialStringTok{f"VPN: S/ }\SpecialCharTok{\{}\NormalTok{vpn}\SpecialCharTok{:,.2f\}}\SpecialStringTok{"}\NormalTok{)}
        
        \ControlFlowTok{if}\NormalTok{ vpn }\OperatorTok{\textgreater{}} \DecValTok{0}\NormalTok{:}
            \BuiltInTok{print}\NormalTok{(}\SpecialStringTok{f"Recomendación: ACEPTAR el proyecto"}\NormalTok{)}
            \BuiltInTok{print}\NormalTok{(}\SpecialStringTok{f"El proyecto genera valor"}\NormalTok{)}
        \ControlFlowTok{elif}\NormalTok{ vpn }\OperatorTok{==} \DecValTok{0}\NormalTok{:}
            \BuiltInTok{print}\NormalTok{(}\SpecialStringTok{f"Recomendación: INDIFERENTE"}\NormalTok{)}
            \BuiltInTok{print}\NormalTok{(}\SpecialStringTok{f"El proyecto no genera ni destruye valor"}\NormalTok{)}
        \ControlFlowTok{else}\NormalTok{:}
            \BuiltInTok{print}\NormalTok{(}\SpecialStringTok{f"Recomendación: RECHAZAR el proyecto"}\NormalTok{)}
            \BuiltInTok{print}\NormalTok{(}\SpecialStringTok{f"El proyecto destruye valor"}\NormalTok{)}
        
        \ControlFlowTok{if}\NormalTok{ periodo\_recuperacion:}
            \BuiltInTok{print}\NormalTok{(}\SpecialStringTok{f"}\CharTok{\textbackslash{}n}\SpecialStringTok{Período de recuperación: }\SpecialCharTok{\{}\NormalTok{periodo\_recuperacion}\SpecialCharTok{\}}\SpecialStringTok{ años"}\NormalTok{)}
        \ControlFlowTok{else}\NormalTok{:}
            \BuiltInTok{print}\NormalTok{(}\SpecialStringTok{f"}\CharTok{\textbackslash{}n}\SpecialStringTok{El proyecto no recupera la inversión en el período analizado"}\NormalTok{)}
        
        \BuiltInTok{print}\NormalTok{(}\SpecialStringTok{f"}\SpecialCharTok{\{}\StringTok{\textquotesingle{}=\textquotesingle{}}\OperatorTok{*}\DecValTok{60}\SpecialCharTok{\}}\SpecialStringTok{"}\NormalTok{)}
        
        \ControlFlowTok{return}\NormalTok{ vpn}
        
    \ControlFlowTok{except} \PreprocessorTok{KeyError} \ImportTok{as}\NormalTok{ e:}
        \BuiltInTok{print}\NormalTok{(}\SpecialStringTok{f"Error: Falta el campo }\SpecialCharTok{\{}\NormalTok{e}\SpecialCharTok{\}}\SpecialStringTok{"}\NormalTok{)}
        \ControlFlowTok{return} \VariableTok{None}
    \ControlFlowTok{except} \PreprocessorTok{Exception} \ImportTok{as}\NormalTok{ e:}
        \BuiltInTok{print}\NormalTok{(}\SpecialStringTok{f"Error en la evaluación: }\SpecialCharTok{\{}\NormalTok{e}\SpecialCharTok{\}}\SpecialStringTok{"}\NormalTok{)}
        \ControlFlowTok{return} \VariableTok{None}

\CommentTok{\# Ejemplo de uso}
\NormalTok{proyecto\_a }\OperatorTok{=}\NormalTok{ \{}
    \StringTok{\textquotesingle{}nombre\textquotesingle{}}\NormalTok{: }\StringTok{\textquotesingle{}Ampliación de planta\textquotesingle{}}\NormalTok{,}
    \StringTok{\textquotesingle{}inversion\_inicial\textquotesingle{}}\NormalTok{: }\DecValTok{500000}\NormalTok{,}
    \StringTok{\textquotesingle{}flujos\_anuales\textquotesingle{}}\NormalTok{: [}\DecValTok{150000}\NormalTok{, }\DecValTok{180000}\NormalTok{, }\DecValTok{200000}\NormalTok{, }\DecValTok{220000}\NormalTok{, }\DecValTok{200000}\NormalTok{],}
    \StringTok{\textquotesingle{}tasa\_descuento\textquotesingle{}}\NormalTok{: }\FloatTok{0.12}
\NormalTok{\}}

\NormalTok{vpn\_a }\OperatorTok{=}\NormalTok{ evaluar\_proyecto(proyecto\_a)}

\BuiltInTok{print}\NormalTok{(}\StringTok{"}\CharTok{\textbackslash{}n}\StringTok{"} \OperatorTok{+} \StringTok{"="}\OperatorTok{*}\DecValTok{80} \OperatorTok{+} \StringTok{"}\CharTok{\textbackslash{}n}\StringTok{"}\NormalTok{)}

\NormalTok{proyecto\_b }\OperatorTok{=}\NormalTok{ \{}
    \StringTok{\textquotesingle{}nombre\textquotesingle{}}\NormalTok{: }\StringTok{\textquotesingle{}Nueva línea de productos\textquotesingle{}}\NormalTok{,}
    \StringTok{\textquotesingle{}inversion\_inicial\textquotesingle{}}\NormalTok{: }\DecValTok{800000}\NormalTok{,}
    \StringTok{\textquotesingle{}flujos\_anuales\textquotesingle{}}\NormalTok{: [}\DecValTok{100000}\NormalTok{, }\DecValTok{150000}\NormalTok{, }\DecValTok{200000}\NormalTok{, }\DecValTok{250000}\NormalTok{, }\DecValTok{300000}\NormalTok{],}
    \StringTok{\textquotesingle{}tasa\_descuento\textquotesingle{}}\NormalTok{: }\FloatTok{0.12}
\NormalTok{\}}

\NormalTok{vpn\_b }\OperatorTok{=}\NormalTok{ evaluar\_proyecto(proyecto\_b)}
\end{Highlighting}
\end{Shaded}

\subsection{Ejercicio 2: Simulador de política
monetaria}\label{ejercicio-2-simulador-de-poluxedtica-monetaria}

\begin{Shaded}
\begin{Highlighting}[]
\KeywordTok{def}\NormalTok{ simular\_politica\_monetaria(estado\_economia):}
    \CommentTok{"""}
\CommentTok{    Simula decisiones de política monetaria según el estado de la economía}
\CommentTok{    """}
    
    \BuiltInTok{print}\NormalTok{(}\StringTok{"SIMULADOR DE POLÍTICA MONETARIA"}\NormalTok{)}
    \BuiltInTok{print}\NormalTok{(}\StringTok{"="} \OperatorTok{*} \DecValTok{60}\NormalTok{)}
    
    \ControlFlowTok{try}\NormalTok{:}
        \CommentTok{\# Extraer indicadores}
\NormalTok{        inflacion\_actual }\OperatorTok{=}\NormalTok{ estado\_economia[}\StringTok{\textquotesingle{}inflacion\_actual\textquotesingle{}}\NormalTok{]}
\NormalTok{        inflacion\_esperada }\OperatorTok{=}\NormalTok{ estado\_economia[}\StringTok{\textquotesingle{}inflacion\_esperada\textquotesingle{}}\NormalTok{]}
\NormalTok{        meta\_inflacion }\OperatorTok{=}\NormalTok{ estado\_economia[}\StringTok{\textquotesingle{}meta\_inflacion\textquotesingle{}}\NormalTok{]}
\NormalTok{        brecha\_producto }\OperatorTok{=}\NormalTok{ estado\_economia[}\StringTok{\textquotesingle{}brecha\_producto\textquotesingle{}}\NormalTok{]}
\NormalTok{        tasa\_actual }\OperatorTok{=}\NormalTok{ estado\_economia[}\StringTok{\textquotesingle{}tasa\_interes\_actual\textquotesingle{}}\NormalTok{]}
        
        \BuiltInTok{print}\NormalTok{(}\SpecialStringTok{f"}\CharTok{\textbackslash{}n}\SpecialStringTok{Indicadores macroeconómicos:"}\NormalTok{)}
        \BuiltInTok{print}\NormalTok{(}\SpecialStringTok{f"  Inflación actual: }\SpecialCharTok{\{}\NormalTok{inflacion\_actual}\SpecialCharTok{:.2f\}}\SpecialStringTok{\%"}\NormalTok{)}
        \BuiltInTok{print}\NormalTok{(}\SpecialStringTok{f"  Inflación esperada: }\SpecialCharTok{\{}\NormalTok{inflacion\_esperada}\SpecialCharTok{:.2f\}}\SpecialStringTok{\%"}\NormalTok{)}
        \BuiltInTok{print}\NormalTok{(}\SpecialStringTok{f"  Meta de inflación: }\SpecialCharTok{\{}\NormalTok{meta\_inflacion}\SpecialCharTok{:.2f\}}\SpecialStringTok{\%"}\NormalTok{)}
        \BuiltInTok{print}\NormalTok{(}\SpecialStringTok{f"  Brecha del producto: }\SpecialCharTok{\{}\NormalTok{brecha\_producto}\SpecialCharTok{:+.2f\}}\SpecialStringTok{\%"}\NormalTok{)}
        \BuiltInTok{print}\NormalTok{(}\SpecialStringTok{f"  Tasa de interés actual: }\SpecialCharTok{\{}\NormalTok{tasa\_actual}\SpecialCharTok{:.2f\}}\SpecialStringTok{\%"}\NormalTok{)}
        
        \CommentTok{\# Regla de Taylor simplificada}
        \CommentTok{\# r = r* + π + 0.5(π {-} π*) + 0.5(y)}
        \CommentTok{\# donde:}
        \CommentTok{\# r = tasa de interés nominal}
        \CommentTok{\# r* = tasa de interés real neutral (asumimos 2\%)}
        \CommentTok{\# π = inflación}
        \CommentTok{\# π* = meta de inflación}
        \CommentTok{\# y = brecha del producto}
        
\NormalTok{        tasa\_neutral }\OperatorTok{=} \FloatTok{2.0}
\NormalTok{        desviacion\_inflacion }\OperatorTok{=}\NormalTok{ inflacion\_esperada }\OperatorTok{{-}}\NormalTok{ meta\_inflacion}
        
        \CommentTok{\# Calcular tasa sugerida según Regla de Taylor}
\NormalTok{        tasa\_sugerida }\OperatorTok{=}\NormalTok{ tasa\_neutral }\OperatorTok{+}\NormalTok{ inflacion\_esperada }\OperatorTok{+} \OperatorTok{\textbackslash{}}
                       \FloatTok{0.5} \OperatorTok{*}\NormalTok{ desviacion\_inflacion }\OperatorTok{+} \OperatorTok{\textbackslash{}}
                       \FloatTok{0.5} \OperatorTok{*}\NormalTok{ brecha\_producto}
        
        \CommentTok{\# Determinar acción}
\NormalTok{        diferencia }\OperatorTok{=}\NormalTok{ tasa\_sugerida }\OperatorTok{{-}}\NormalTok{ tasa\_actual}
        
        \BuiltInTok{print}\NormalTok{(}\SpecialStringTok{f"}\CharTok{\textbackslash{}n}\SpecialCharTok{\{}\StringTok{\textquotesingle{}=\textquotesingle{}}\OperatorTok{*}\DecValTok{60}\SpecialCharTok{\}}\SpecialStringTok{"}\NormalTok{)}
        \BuiltInTok{print}\NormalTok{(}\SpecialStringTok{f"ANÁLISIS:"}\NormalTok{)}
        
        \CommentTok{\# Evaluar presiones inflacionarias}
        \ControlFlowTok{if}\NormalTok{ inflacion\_actual }\OperatorTok{\textgreater{}}\NormalTok{ meta\_inflacion }\OperatorTok{+} \FloatTok{1.0}\NormalTok{:}
            \BuiltInTok{print}\NormalTok{(}\SpecialStringTok{f"⚠️  Inflación significativamente por encima de la meta"}\NormalTok{)}
\NormalTok{            presion\_inflacionaria }\OperatorTok{=} \StringTok{"alta"}
        \ControlFlowTok{elif}\NormalTok{ inflacion\_actual }\OperatorTok{\textgreater{}}\NormalTok{ meta\_inflacion:}
            \BuiltInTok{print}\NormalTok{(}\SpecialStringTok{f"○ Inflación ligeramente por encima de la meta"}\NormalTok{)}
\NormalTok{            presion\_inflacionaria }\OperatorTok{=} \StringTok{"moderada"}
        \ControlFlowTok{elif}\NormalTok{ inflacion\_actual }\OperatorTok{\textless{}}\NormalTok{ meta\_inflacion }\OperatorTok{{-}} \FloatTok{1.0}\NormalTok{:}
            \BuiltInTok{print}\NormalTok{(}\SpecialStringTok{f"○ Inflación significativamente por debajo de la meta"}\NormalTok{)}
\NormalTok{            presion\_inflacionaria }\OperatorTok{=} \StringTok{"baja"}
        \ControlFlowTok{else}\NormalTok{:}
            \BuiltInTok{print}\NormalTok{(}\SpecialStringTok{f"✓ Inflación dentro del rango meta"}\NormalTok{)}
\NormalTok{            presion\_inflacionaria }\OperatorTok{=} \StringTok{"normal"}
        
        \CommentTok{\# Evaluar actividad económica}
        \ControlFlowTok{if}\NormalTok{ brecha\_producto }\OperatorTok{\textgreater{}} \FloatTok{2.0}\NormalTok{:}
            \BuiltInTok{print}\NormalTok{(}\SpecialStringTok{f"⚠️  Economía sobrecalentada"}\NormalTok{)}
\NormalTok{            presion\_demanda }\OperatorTok{=} \StringTok{"alta"}
        \ControlFlowTok{elif}\NormalTok{ brecha\_producto }\OperatorTok{\textgreater{}} \DecValTok{0}\NormalTok{:}
            \BuiltInTok{print}\NormalTok{(}\SpecialStringTok{f"○ Economía por encima de su potencial"}\NormalTok{)}
\NormalTok{            presion\_demanda }\OperatorTok{=} \StringTok{"moderada"}
        \ControlFlowTok{elif}\NormalTok{ brecha\_producto }\OperatorTok{\textless{}} \OperatorTok{{-}}\FloatTok{2.0}\NormalTok{:}
            \BuiltInTok{print}\NormalTok{(}\SpecialStringTok{f"⚠️  Economía en recesión"}\NormalTok{)}
\NormalTok{            presion\_demanda }\OperatorTok{=} \StringTok{"muy baja"}
        \ControlFlowTok{else}\NormalTok{:}
            \BuiltInTok{print}\NormalTok{(}\SpecialStringTok{f"○ Economía por debajo de su potencial"}\NormalTok{)}
\NormalTok{            presion\_demanda }\OperatorTok{=} \StringTok{"baja"}
        
        \CommentTok{\# Recomendación}
        \BuiltInTok{print}\NormalTok{(}\SpecialStringTok{f"}\CharTok{\textbackslash{}n}\SpecialCharTok{\{}\StringTok{\textquotesingle{}=\textquotesingle{}}\OperatorTok{*}\DecValTok{60}\SpecialCharTok{\}}\SpecialStringTok{"}\NormalTok{)}
        \BuiltInTok{print}\NormalTok{(}\SpecialStringTok{f"RECOMENDACIÓN DE POLÍTICA:"}\NormalTok{)}
        
        \ControlFlowTok{if} \BuiltInTok{abs}\NormalTok{(diferencia) }\OperatorTok{\textless{}} \FloatTok{0.25}\NormalTok{:}
\NormalTok{            decision }\OperatorTok{=} \StringTok{"MANTENER"}
\NormalTok{            tasa\_nueva }\OperatorTok{=}\NormalTok{ tasa\_actual}
            \BuiltInTok{print}\NormalTok{(}\SpecialStringTok{f"✓ MANTENER la tasa de interés en }\SpecialCharTok{\{}\NormalTok{tasa\_actual}\SpecialCharTok{:.2f\}}\SpecialStringTok{\%"}\NormalTok{)}
            \BuiltInTok{print}\NormalTok{(}\SpecialStringTok{f"  La tasa actual es adecuada para las condiciones macroeconómicas"}\NormalTok{)}
        
        \ControlFlowTok{elif}\NormalTok{ diferencia }\OperatorTok{\textgreater{}} \DecValTok{0}\NormalTok{:}
            \ControlFlowTok{if} \BuiltInTok{abs}\NormalTok{(diferencia) }\OperatorTok{\textgreater{}} \FloatTok{0.75}\NormalTok{:}
\NormalTok{                ajuste }\OperatorTok{=} \FloatTok{0.50}
\NormalTok{                decision }\OperatorTok{=} \StringTok{"SUBIR FUERTE"}
            \ControlFlowTok{else}\NormalTok{:}
\NormalTok{                ajuste }\OperatorTok{=} \FloatTok{0.25}
\NormalTok{                decision }\OperatorTok{=} \StringTok{"SUBIR"}
            
\NormalTok{            tasa\_nueva }\OperatorTok{=}\NormalTok{ tasa\_actual }\OperatorTok{+}\NormalTok{ ajuste}
            \BuiltInTok{print}\NormalTok{(}\SpecialStringTok{f"↑ }\SpecialCharTok{\{}\NormalTok{decision}\SpecialCharTok{\}}\SpecialStringTok{ la tasa de interés a }\SpecialCharTok{\{}\NormalTok{tasa\_nueva}\SpecialCharTok{:.2f\}}\SpecialStringTok{\%"}\NormalTok{)}
            \BuiltInTok{print}\NormalTok{(}\SpecialStringTok{f"  Ajuste: +}\SpecialCharTok{\{}\NormalTok{ajuste}\SpecialCharTok{:.2f\}}\SpecialStringTok{ puntos porcentuales"}\NormalTok{)}
            
            \ControlFlowTok{if}\NormalTok{ presion\_inflacionaria }\KeywordTok{in}\NormalTok{ [}\StringTok{"alta"}\NormalTok{, }\StringTok{"moderada"}\NormalTok{]:}
                \BuiltInTok{print}\NormalTok{(}\SpecialStringTok{f"  Justificación: Controlar presiones inflacionarias"}\NormalTok{)}
            \ControlFlowTok{if}\NormalTok{ presion\_demanda }\OperatorTok{==} \StringTok{"alta"}\NormalTok{:}
                \BuiltInTok{print}\NormalTok{(}\SpecialStringTok{f"  Justificación: Enfriar la economía sobrecalentada"}\NormalTok{)}
        
        \ControlFlowTok{else}\NormalTok{:}
            \ControlFlowTok{if} \BuiltInTok{abs}\NormalTok{(diferencia) }\OperatorTok{\textgreater{}} \FloatTok{0.75}\NormalTok{:}
\NormalTok{                ajuste }\OperatorTok{=} \OperatorTok{{-}}\FloatTok{0.50}
\NormalTok{                decision }\OperatorTok{=} \StringTok{"BAJAR FUERTE"}
            \ControlFlowTok{else}\NormalTok{:}
\NormalTok{                ajuste }\OperatorTok{=} \OperatorTok{{-}}\FloatTok{0.25}
\NormalTok{                decision }\OperatorTok{=} \StringTok{"BAJAR"}
            
\NormalTok{            tasa\_nueva }\OperatorTok{=}\NormalTok{ tasa\_actual }\OperatorTok{+}\NormalTok{ ajuste}
            \BuiltInTok{print}\NormalTok{(}\SpecialStringTok{f"↓ }\SpecialCharTok{\{}\NormalTok{decision}\SpecialCharTok{\}}\SpecialStringTok{ la tasa de interés a }\SpecialCharTok{\{}\NormalTok{tasa\_nueva}\SpecialCharTok{:.2f\}}\SpecialStringTok{\%"}\NormalTok{)}
            \BuiltInTok{print}\NormalTok{(}\SpecialStringTok{f"  Ajuste: }\SpecialCharTok{\{}\NormalTok{ajuste}\SpecialCharTok{:.2f\}}\SpecialStringTok{ puntos porcentuales"}\NormalTok{)}
            
            \ControlFlowTok{if}\NormalTok{ presion\_inflacionaria }\OperatorTok{==} \StringTok{"baja"}\NormalTok{:}
                \BuiltInTok{print}\NormalTok{(}\SpecialStringTok{f"  Justificación: Evitar presiones deflacionarias"}\NormalTok{)}
            \ControlFlowTok{if}\NormalTok{ presion\_demanda }\KeywordTok{in}\NormalTok{ [}\StringTok{"baja"}\NormalTok{, }\StringTok{"muy baja"}\NormalTok{]:}
                \BuiltInTok{print}\NormalTok{(}\SpecialStringTok{f"  Justificación: Estimular la actividad económica"}\NormalTok{)}
        
        \CommentTok{\# Efectos esperados}
        \BuiltInTok{print}\NormalTok{(}\SpecialStringTok{f"}\CharTok{\textbackslash{}n}\SpecialStringTok{Efectos esperados:"}\NormalTok{)}
        \ControlFlowTok{if}\NormalTok{ decision }\KeywordTok{in}\NormalTok{ [}\StringTok{"SUBIR"}\NormalTok{, }\StringTok{"SUBIR FUERTE"}\NormalTok{]:}
            \BuiltInTok{print}\NormalTok{(}\SpecialStringTok{f"  • Encarecimiento del crédito"}\NormalTok{)}
            \BuiltInTok{print}\NormalTok{(}\SpecialStringTok{f"  • Reducción del consumo e inversión"}\NormalTok{)}
            \BuiltInTok{print}\NormalTok{(}\SpecialStringTok{f"  • Desaceleración de la inflación"}\NormalTok{)}
            \BuiltInTok{print}\NormalTok{(}\SpecialStringTok{f"  • Apreciación de la moneda local"}\NormalTok{)}
        \ControlFlowTok{elif}\NormalTok{ decision }\KeywordTok{in}\NormalTok{ [}\StringTok{"BAJAR"}\NormalTok{, }\StringTok{"BAJAR FUERTE"}\NormalTok{]:}
            \BuiltInTok{print}\NormalTok{(}\SpecialStringTok{f"  • Abaratamiento del crédito"}\NormalTok{)}
            \BuiltInTok{print}\NormalTok{(}\SpecialStringTok{f"  • Estímulo al consumo e inversión"}\NormalTok{)}
            \BuiltInTok{print}\NormalTok{(}\SpecialStringTok{f"  • Impulso a la actividad económica"}\NormalTok{)}
            \BuiltInTok{print}\NormalTok{(}\SpecialStringTok{f"  • Depreciación de la moneda local"}\NormalTok{)}
        \ControlFlowTok{else}\NormalTok{:}
            \BuiltInTok{print}\NormalTok{(}\SpecialStringTok{f"  • Mantenimiento de las condiciones actuales"}\NormalTok{)}
            \BuiltInTok{print}\NormalTok{(}\SpecialStringTok{f"  • Monitoreo continuo de indicadores"}\NormalTok{)}
        
        \BuiltInTok{print}\NormalTok{(}\SpecialStringTok{f"}\SpecialCharTok{\{}\StringTok{\textquotesingle{}=\textquotesingle{}}\OperatorTok{*}\DecValTok{60}\SpecialCharTok{\}}\SpecialStringTok{"}\NormalTok{)}
        
        \ControlFlowTok{return}\NormalTok{ \{}
            \StringTok{\textquotesingle{}decision\textquotesingle{}}\NormalTok{: decision,}
            \StringTok{\textquotesingle{}tasa\_nueva\textquotesingle{}}\NormalTok{: tasa\_nueva,}
            \StringTok{\textquotesingle{}tasa\_sugerida\_taylor\textquotesingle{}}\NormalTok{: tasa\_sugerida}
\NormalTok{        \}}
        
    \ControlFlowTok{except} \PreprocessorTok{Exception} \ImportTok{as}\NormalTok{ e:}
        \BuiltInTok{print}\NormalTok{(}\SpecialStringTok{f"Error en la simulación: }\SpecialCharTok{\{}\NormalTok{e}\SpecialCharTok{\}}\SpecialStringTok{"}\NormalTok{)}
        \ControlFlowTok{return} \VariableTok{None}

\CommentTok{\# Escenario 1: Inflación alta}
\BuiltInTok{print}\NormalTok{(}\StringTok{"ESCENARIO 1: Presiones inflacionarias"}\NormalTok{)}
\NormalTok{escenario\_1 }\OperatorTok{=}\NormalTok{ \{}
    \StringTok{\textquotesingle{}inflacion\_actual\textquotesingle{}}\NormalTok{: }\FloatTok{5.2}\NormalTok{,}
    \StringTok{\textquotesingle{}inflacion\_esperada\textquotesingle{}}\NormalTok{: }\FloatTok{5.5}\NormalTok{,}
    \StringTok{\textquotesingle{}meta\_inflacion\textquotesingle{}}\NormalTok{: }\FloatTok{3.0}\NormalTok{,}
    \StringTok{\textquotesingle{}brecha\_producto\textquotesingle{}}\NormalTok{: }\FloatTok{1.5}\NormalTok{,}
    \StringTok{\textquotesingle{}tasa\_interes\_actual\textquotesingle{}}\NormalTok{: }\FloatTok{5.5}
\NormalTok{\}}
\NormalTok{resultado\_1 }\OperatorTok{=}\NormalTok{ simular\_politica\_monetaria(escenario\_1)}

\BuiltInTok{print}\NormalTok{(}\StringTok{"}\CharTok{\textbackslash{}n}\StringTok{"} \OperatorTok{+} \StringTok{"="}\OperatorTok{*}\DecValTok{80} \OperatorTok{+} \StringTok{"}\CharTok{\textbackslash{}n}\StringTok{"}\NormalTok{)}

\CommentTok{\# Escenario 2: Recesión}
\BuiltInTok{print}\NormalTok{(}\StringTok{"ESCENARIO 2: Recesión económica"}\NormalTok{)}
\NormalTok{escenario\_2 }\OperatorTok{=}\NormalTok{ \{}
    \StringTok{\textquotesingle{}inflacion\_actual\textquotesingle{}}\NormalTok{: }\FloatTok{2.0}\NormalTok{,}
    \StringTok{\textquotesingle{}inflacion\_esperada\textquotesingle{}}\NormalTok{: }\FloatTok{1.8}\NormalTok{,}
    \StringTok{\textquotesingle{}meta\_inflacion\textquotesingle{}}\NormalTok{: }\FloatTok{3.0}\NormalTok{,}
    \StringTok{\textquotesingle{}brecha\_producto\textquotesingle{}}\NormalTok{: }\OperatorTok{{-}}\FloatTok{3.5}\NormalTok{,}
    \StringTok{\textquotesingle{}tasa\_interes\_actual\textquotesingle{}}\NormalTok{: }\FloatTok{4.0}
\NormalTok{\}}
\NormalTok{resultado\_2 }\OperatorTok{=}\NormalTok{ simular\_politica\_monetaria(escenario\_2)}
\end{Highlighting}
\end{Shaded}

\section{Conclusión}\label{conclusiuxf3n}

En esta guía hemos explorado la ejecución condicional en Python:

\textbf{Expresiones booleanas}

Expresiones que evalúan a True o False. Son la base de la toma de
decisiones en programación.

\textbf{Operadores de comparación}

Permiten comparar valores: \texttt{==}, \texttt{!=},
\texttt{\textgreater{}}, \texttt{\textless{}}, \texttt{\textgreater{}=},
\texttt{\textless{}=}.

\textbf{Operadores lógicos}

Combinan expresiones booleanas: \texttt{and}, \texttt{or}, \texttt{not}.
La precedencia es: \texttt{not} \textgreater{} \texttt{and}
\textgreater{} \texttt{or}.

\textbf{Estructuras condicionales}

\texttt{if} para ejecución simple, \texttt{if-else} para alternativas,
\texttt{if-elif-else} para múltiples condiciones.

\textbf{Condicionales anidados}

Permiten lógica compleja pero deben usarse con moderación para mantener
legibilidad.

\textbf{Manejo de excepciones}

\texttt{try-except} permite capturar y manejar errores sin detener el
programa.

\textbf{Evaluación de cortocircuito}

Python evalúa expresiones booleanas eficientemente, deteniéndose cuando
el resultado es determinado.

\section{Próximos pasos}\label{pruxf3ximos-pasos}

En la siguiente guía (Guía 5: Iteraciones) exploraremos:

\begin{itemize}
\tightlist
\item
  Actualización de variables
\item
  Bucles definidos con for
\item
  Iteración doble, múltiple y anidada
\item
  Bucles while
\item
  List comprehension
\end{itemize}

Las iteraciones nos permitirán procesar grandes cantidades de datos
económicos de manera eficiente y automatizada.

\section{Recursos adicionales}\label{recursos-adicionales}

Para profundizar en ejecución condicional:

\begin{itemize}
\tightlist
\item
  Python Control Flow: docs.python.org/tutorial/controlflow.html
\item
  Práctica con casos económicos reales
\item
  Estudio de algoritmos de decisión en economía
\end{itemize}

\section{Publicaciones Similares}\label{publicaciones-similares}

Si te interesó este artículo, te recomendamos que explores otros blogs y
recursos relacionados que pueden ampliar tus conocimientos. Aquí te dejo
algunas sugerencias:

\begin{enumerate}
\def\labelenumi{\arabic{enumi}.}
\tightlist
\item
  \href{https://numerus-scriptum.netlify.app/python/2020-06-19-instalacion-de-anaconda/index.pdf}{\faIcon{file-pdf}}
  \href{https://numerus-scriptum.netlify.app/python/2020-06-19-instalacion-de-anaconda}{Instalacion
  De Anaconda}
\item
  \href{https://numerus-scriptum.netlify.app/python/2020-06-20-configurar-entorno-virtual-python-anaconda/index.pdf}{\faIcon{file-pdf}}
  \href{https://numerus-scriptum.netlify.app/python/2020-06-20-configurar-entorno-virtual-python-anaconda}{Configurar
  Entorno Virtual Python Anaconda}
\item
  \href{https://numerus-scriptum.netlify.app/python/2021-04-17-01-introducion-a-la-programacion-con-python/index.pdf}{\faIcon{file-pdf}}
  \href{https://numerus-scriptum.netlify.app/python/2021-04-17-01-introducion-a-la-programacion-con-python}{01
  Introducion A La Programacion Con Python}
\item
  \href{https://numerus-scriptum.netlify.app/python/2021-05-31-02-variables-expresiones-y-statements-con-python/index.pdf}{\faIcon{file-pdf}}
  \href{https://numerus-scriptum.netlify.app/python/2021-05-31-02-variables-expresiones-y-statements-con-python}{02
  Variables Expresiones Y Statements Con Python}
\item
  \href{https://numerus-scriptum.netlify.app/python/2021-06-07-03-objetos-de-python/index.pdf}{\faIcon{file-pdf}}
  \href{https://numerus-scriptum.netlify.app/python/2021-06-07-03-objetos-de-python}{03
  Objetos De Python}
\item
  \href{https://numerus-scriptum.netlify.app/python/2021-06-14-04-ejecucion-condicional-con-python/index.pdf}{\faIcon{file-pdf}}
  \href{https://numerus-scriptum.netlify.app/python/2021-06-14-04-ejecucion-condicional-con-python}{04
  Ejecucion Condicional Con Python}
\item
  \href{https://numerus-scriptum.netlify.app/python/2021-06-21-05-iteraciones-con-python/index.pdf}{\faIcon{file-pdf}}
  \href{https://numerus-scriptum.netlify.app/python/2021-06-21-05-iteraciones-con-python}{05
  Iteraciones Con Python}
\item
  \href{https://numerus-scriptum.netlify.app/python/2021-08-16-06-funciones-con-python/index.pdf}{\faIcon{file-pdf}}
  \href{https://numerus-scriptum.netlify.app/python/2021-08-16-06-funciones-con-python}{06
  Funciones Con Python}
\item
  \href{https://numerus-scriptum.netlify.app/python/2021-08-23-07-dataframes-con-python/index.pdf}{\faIcon{file-pdf}}
  \href{https://numerus-scriptum.netlify.app/python/2021-08-23-07-dataframes-con-python}{07
  Dataframes Con Python}
\item
  \href{https://numerus-scriptum.netlify.app/python/2021-11-29-08-prediccion-y-metrica-de-performance-con-python/index.pdf}{\faIcon{file-pdf}}
  \href{https://numerus-scriptum.netlify.app/python/2021-11-29-08-prediccion-y-metrica-de-performance-con-python}{08
  Prediccion Y Metrica De Performance Con Python}
\item
  \href{https://numerus-scriptum.netlify.app/python/2021-12-06-09-metodos-de-machine-learning-para-clasificacion-con-python/index.pdf}{\faIcon{file-pdf}}
  \href{https://numerus-scriptum.netlify.app/python/2021-12-06-09-metodos-de-machine-learning-para-clasificacion-con-python}{09
  Metodos De Machine Learning Para Clasificacion Con Python}
\item
  \href{https://numerus-scriptum.netlify.app/python/2021-12-13-10-metodos-de-machine-learning-para-regresion-con-python/index.pdf}{\faIcon{file-pdf}}
  \href{https://numerus-scriptum.netlify.app/python/2021-12-13-10-metodos-de-machine-learning-para-regresion-con-python}{10
  Metodos De Machine Learning Para Regresion Con Python}
\item
  \href{https://numerus-scriptum.netlify.app/python/2022-10-31-11-validacion-cruzada-y-composicion-del-modelo-con-python/index.pdf}{\faIcon{file-pdf}}
  \href{https://numerus-scriptum.netlify.app/python/2022-10-31-11-validacion-cruzada-y-composicion-del-modelo-con-python}{11
  Validacion Cruzada Y Composicion Del Modelo Con Python}
\item
  \href{https://numerus-scriptum.netlify.app/python/2025-05-10-visualizacion-de-datos-con-python/index.pdf}{\faIcon{file-pdf}}
  \href{https://numerus-scriptum.netlify.app/python/2025-05-10-visualizacion-de-datos-con-python}{Visualizacion
  De Datos Con Python}
\end{enumerate}

Esperamos que encuentres estas publicaciones igualmente interesantes y
útiles. ¡Disfruta de la lectura!






\end{document}
